\documentclass[italian,aspectratio=169]{beamer}
\usepackage[italian]{babel}
\usepackage[utf8]{inputenc}
\usepackage[T1]{fontenc}
\usepackage{lmodern}
\usepackage{amsmath,amssymb}
\usepackage{mathtools}
\usepackage{textcomp}
\usepackage{booktabs}
\usepackage{hyperref}
\usepackage{xcolor}
\usepackage{listings}
% Metropolis theme
\usepackage{beamerthememetropolis}
\metroset{
  progressbar=frametitle,
  sectionpage=progressbar,
  numbering=counter,
  block=fill,
}

% Color palette
\definecolor{cb}{RGB}{41,128,185}   % Blue  - Intuizione
\definecolor{co}{RGB}{231,76,60}    % Red   - All'orale
\definecolor{cg}{RGB}{39,174,96}    % Green - Spiegazioni
\definecolor{cp}{RGB}{142,68,173}   % Purple - Termini

\setbeamercolor{background canvas}{bg=white!98}
\setbeamercolor{alerted text}{fg=co}
\setbeamercolor{example text}{fg=cg}
\setbeamercolor{normal text}{fg=black!90}

% Custom blocks (Metropolis-compatible, with fill)
\newenvironment{intuizione}{%
  \setbeamercolor{block title}{fg=white,bg=cb}%
  \setbeamercolor{block body}{fg=black,bg=cb!5!white}%
  \begin{block}{Intuizione}}{\end{block}}
\newenvironment{orale}{%
  \setbeamercolor{block title}{fg=white,bg=co}%
  \setbeamercolor{block body}{fg=black,bg=co!5!white}%
  \begin{block}{All'orale}}{\end{block}}
\newenvironment{spiegazione}{%
  \setbeamercolor{block title}{fg=white,bg=cg}%
  \setbeamercolor{block body}{fg=black,bg=cg!5!white}%
  \begin{block}{Spiegazioni}}{\end{block}}
\newenvironment{terminiblock}{%
  \setbeamercolor{block title}{fg=white,bg=cp}%
  \setbeamercolor{block body}{fg=black,bg=cp!5!white}%
  \begin{block}{Termini}}{\end{block}}

% Compact code blocks with shadow
\lstdefinestyle{pp}{
  language=Python,
  basicstyle=\ttfamily\small,
  backgroundcolor=\color{gray!3},
  frame=shadowbox,
  rulesepcolor=\color{gray!30},
  breaklines=true,
  showstringspaces=false,
  commentstyle=\color{gray},
  keywordstyle=\color{cb},
}

% Slightly smaller text for all frames (better density)
\apptocmd{\frame}{\small}{}{}

% Section pages (Metropolis built-in handles the rest)
\title{Computational Imaging}
\subtitle{Studio per l'Esame Orale}
\date{}

\begin{document}
\frame{\titlepage}
\frame{\frametitle{Indice}\tableofcontents}

\begin{frame}{}
\textbf{Unibo - Laurea Magistrale in Informatica}: 
Modulo 1: Prof.ssa Elena Loli Piccolomini | Modulo 2: Davide Evangelista
Link: [Modulo 1](https://elenaloli.github.io/computational-imaging/intro.html) | [Modulo 2](https://devangelista2.github.io/computational-imaging/2025-26/intro/environment-setup.html)
GitHub Modulo 1: [elenaloli/computational-imaging](https://github.com/elenaloli/computational-imaging) | GitHub Modulo 2: [devangelista2/computational-imaging](https://github.com/devangelista2/computational-imaging)
\medskip\noindent\rule{\textwidth}{0.4pt}\medskip
\end{frame}

\section{Indice}
\begin{frame}{Indice}
\tableofcontents[currentsection]
\end{frame}

\begin{frame}[allowframebreaks]{}
\begin{itemize}
\item \textbf{PARTE I - Fondamenti} (Modulo 1, sez. 1-6) `[p.1-11]`
\item 1. Introduzione `[p.1-2]`
\item 2. Pixel Processing `[p.2-3]`
\item 3. Filtri e Convoluzione `[p.4-5]`
\item 4. Trasformata di Fourier `[p.6-8]`
\item 5. Problemi Inversi `[p.8-10]`
\item 6. Regolarizzazione `[p.10-11]`
\item \textbf{PARTE II - Deep Learning} (Modulo 2, sez. 7-16) `[p.12-31]`
\item 7. Processing per Reti Neurali `[p.12-13]`
\item 8. PyTorch Essentials `[p.14-15]`
\item 9. ML -> Neural Networks `[p.15-16]`
\item 10. CNN `[p.17-18]`
\item 11. UNet `[p.18-20]`
\item 12. ViT \& Loss Design `[p.20-22]`
\item 13. Cross-Domain `[p.22-23]`
\item 14. VAE \& GAN `[p.23-26]`
\item 15. Diffusion Models `[p.26-28]`
\item 16. Diffusion per Problemi Inversi `[p.29-31]`
\item \textbf{Schede Riassuntive} `[p.32-34]`
\item \textbf{Domande d'Esame} `[p.35-42]`
\end{itemize}
\medskip\noindent\rule{\textwidth}{0.4pt}\medskip
\end{frame}

\section{PARTE I - FONDAMENTI DI IMAGING (Modulo 1)}
\begin{frame}{PARTE I - FONDAMENTI DI IMAGING (Modulo 1)}
\tableofcontents[currentsection]
\end{frame}

\begin{frame}{}
\medskip\noindent\rule{\textwidth}{0.4pt}\medskip
\end{frame}

\begin{frame}[allowframebreaks]{1. Introduzione al Computational Imaging}
\begin{terminiblock}discretizzazione = convertire segnale continuo in discreto · Nyquist = $f_s \ge 2f_{max}$ · aliasing = artefatto da sottocampionamento (ruote carro) · DICOM = formato medico con metadati · DCT = base JPEG separa frequenze · pixel = unità minima\end{terminiblock}
\textit{capire come campionare senza perdere informazione e comprimere immagini. Dalla fotografia alla risonanza magnetica, senza Nyquist avremmo solo artefatti.}
\textbf{Esempio}: fotografi un mulino a vento a bassa frequenza: le pale sembrano ferme = aliasing. JPEG comprime scartando le frequenze che l'occhio non vede.
\begin{spiegazione}Il \textbf{Computational Imaging (CI)} ruota attorno a un'idea semplice ma potente: invertire il \textbf{forward model} $\boldsymbol{y} = A\boldsymbol{x} + \boldsymbol{e}$. Misuri dati indiretti e rumorosi ($\boldsymbol{y}$, es. il sinogramma grezzo di una TAC) e da lì ricostruisci l'immagine incognita ($\boldsymbol{x}$, es. la sezione nitida del corpo). È come capire cosa c'è in una scatola chiusa dal rumore che fa scuotendola: non vedi il contenuto, ma dalle vibrazioni deduci forma e materiale.\end{spiegazione}

Il mondo reale è continuo (onde, densità dei tessuti), ma i computer lavorano con dati discreti. Qui entra il \textbf{teorema di Nyquist-Shannon}: per campionare senza perdita serve $f_s \ge 2f_{max}$. Sotto questa soglia le alte frequenze si \textbf{ripiegano} su quelle basse — è l'\textbf{aliasing}, l'effetto "ruote del carro che girano all'indietro" nei film western (24 fps insufficienti). La frequenza critica $f_s/2$ si chiama appunto \textbf{frequenza di Nyquist}.

Per la compressione, il JPEG usa la \textbf{DCT} (Discrete Cosine Transform), che separa basse frequenze (zone uniformi, un cielo azzurro) dalle alte (dettagli fini, texture di capelli). Poi \textbf{quantizza} brutalmente le alte: l'occhio umano non le percepisce bene, quindi buttarle via è quasi invisibile. È come descrivere una spiaggia con «sabbia colorata» invece di elencare ogni granello. Il rapporto di compressione tipico è 10:1-20:1.

In medicina lo standard è \textbf{DICOM}: non solo pixel, ma metadati completi del paziente e dello scanner — fondamentali per la diagnosi. In Python: OpenCV carica in BGR (uint8), Matplotlib in RGB (float [0,1]), imageio supporta anche DICOM.
\begin{orale}«Cos'è il CI?» = invertire $y=Ax+e$ per ricostruire immagini da misure indirette (CT, risonanza). Nyquist: se campioni a meno di $2f_{max}$ le frequenze si ripiegano (aliasing). Ruote carro girano all'indietro. Domanda tipica: «Perché serve Nyquist?» = per evitare aliasing.\end{orale}
\medskip\noindent\rule{\textwidth}{0.4pt}\medskip
\textbf{Computational Imaging (CI)} = campo interdisciplinare che unisce acquisizione, elaborazione e analisi delle immagini attraverso modelli computazionali.
Tre pilastri:
\begin{itemize}
\item \textbf{Image Processing}: elaborare un'immagine per migliorarla o estrarre informazione (filtraggio, enhancement, restauro)
\item \textbf{Computer Vision}: estrarre informazione semantica dalle immagini (riconoscimento, segmentazione, ricostruzione 3D)
\item \textbf{Computer Graphics}: sintetizzare immagini a partire da modelli (rendering, simulazione)
\end{itemize}
\textbf{Immagine Digitale}
Un'immagine digitale e una matrice discreta $I \in \mathbb{R}^{M \times N}$ (o $\mathbb{R}^{M \times N \times C}$ per immagini a colori).
\begin{itemize}
\item \textbf{Risoluzione spaziale}: $M \times N$ pixel
\item \textbf{Profondita di bit}: numero di valori per pixel (8-bit = 256 livelli)
\item \textbf{Canali}: 1 (grigio), 3 (RGB), 4 (RGBA)
\end{itemize}
\textbf{Formati file}:
\begin{itemize}
\item \textbf{DICOM} (Digital Imaging and Communications in Medicine): standard per imaging medico, include metadati clinici
\item \textbf{PNG}: compressione lossless
\item \textbf{JPEG}: compressione lossy (DCT + quantizzazione)
\end{itemize}
\textbf{Lettura immagini in Python}:
\end{frame}

\begin{lstlisting}[style=pp]
import cv2          # OpenCV: BGR, uint8, shape (H,W,C)
import matplotlib   # RGB, float [0,1]
from PIL import Image  # Pillow: RGB, vari formati
import imageio      # versatile, supporta DICOM
\end{lstlisting}
\begin{frame}{}
\medskip\noindent\rule{\textwidth}{0.4pt}\medskip
\end{frame}

\begin{frame}[allowframebreaks]{2. Pixel Processing}
\begin{terminiblock}istogramma = conteggio frequenza livelli grigio · CDF = somma cumulativa · PDF = istogramma normalizzato · equalizzazione = CDF come lookup table · CLAHE = equalizzazione locale a blocchi con clipping · gaussiano = rumore normale · salt\&pepper = pixel bianchi/neri · Poisson = rumore dipendente dal segnale\end{terminiblock}
\textit{migliorare la qualita di un'immagine grezza. Aumentare contrasto, regolare luminosita, correggere colori — operazioni base prima di qualsiasi analisi.}
\textbf{Esempio}: foto in controluce: il viso e tutto nero. Equalizzazione spalma l'istogramma e ora vedi occhi e naso.
\begin{spiegazione}Il \textbf{pixel processing} opera punto-a-punto, senza guardare i vicini. Il punto di partenza è l'\textbf{istogramma} $h(k)$: conta quanti pixel hanno una certa intensità $k$ — è un censimento: per ogni "età" (livello di grigio 0-255), quante "persone" (pixel) ci sono? Normalizzando ottieni la \textbf{PDF} $p(k) = h(k)/(M\cdot N)$, che stima la probabilità di incontrare un dato livello. Istogramma tutto concentrato in poche barre? Poco contrasto. Distribuito su tutto 0-255? Buon contrasto.\end{spiegazione}

L'\textbf{equalizzazione} prende la \textbf{CDF} $T(k) = \sum_{i=0}^k p(i)$ e la usa come tabella di lookup: ogni pixel di intensità $k$ viene rimpiazzato da $T(k) \times 255$. Il risultato è un istogramma il più uniforme possibile — le ombre si aprono, le luci si comprimono, il contrasto globale esplode. È come il cursore "ombre/luci" di Photoshop. \textbf{Limite}: essendo globale, amplifica anche il rumore nelle zone piatte (dove non ci sono dettagli, solo grana).

Qui entra il \textbf{CLAHE} (Contrast Limited Adaptive Histogram Equalization). Divide l'immagine in blocchi (es. 8×8) ed equalizza ciascun blocco separatamente → preserva il contrasto locale. Ma c'è un secondo trucco: il \textbf{contrast clipping} — taglia i picchi dell'istogramma di ogni blocco PRIMA di equalizzare. Così una singola intensità dominante (es. una zona di cielo uniforme) non viene stirata all'estremo, evitando di amplificare il rumore.

Sul fronte \textbf{rumore}, ci sono quattro modelli chiave. Il \textbf{gaussiano} $\mathcal{N}(0,\sigma^2)$ è additivo e indipendente dal segnale: grana fine e uniforme, classica dei sensori CMOS a ISO alti. Il \textbf{salt \& pepper} è impulsivo: pixel casuali schizzano a 0 (nero) o 255 (bianco) — sensore danneggiato o errore di trasmissione. Il \textbf{Poisson} (shot noise) dipende dall'intensità: più luce = più fotoni = più rumore in valore assoluto, ma il rapporto segnale-rumore migliora. Lo \textbf{speckle} è moltiplicativo, tipico di radar e ultrasuoni.

Il \textbf{filtro mediano} merita un'attenzione speciale: NON è lineare, NON è una convoluzione. Prende tutti i pixel nella finestra $k \times k$, li ordina, e sceglie la \textbf{mediana} (il valore centrale). Perché è geniale contro il salt \& pepper? Un pixel impulsivo bianco (255) in mezzo a pixel scuri (10-15), dopo l'ordinamento, finisce all'estremità — la mediana lo ignora completamente. Inoltre \textbf{preserva i bordi}: sul confine tra due regioni, la mediana sceglie il valore del lato con più pixel nella finestra, invece di fare una media che sfocherebbe tutto. Costo: $O(k^2 \log k)$ per l'ordinamento (vs il gaussiano che è $O(k^2)$).
\begin{orale}Equalizzazione = CDF come lookup table. CLAHE divide in blocchi + taglia picchi di contrasto (non amplifica rumore). Se mostrano istogramma stretto: «poco contrasto, equalizzazione allarga». Domanda: «equalizzazione vs CLAHE?» = CLAHE locale + evita amplificazione rumore.\end{orale}
\medskip\noindent\rule{\textwidth}{0.4pt}\medskip
Operazioni che agiscono su ogni pixel \textbf{indipendentemente} dagli altri (punto-a-punto).
\textbf{2.1 Istogramma}
L'istogramma $h(k)$ conta il numero di pixel con intensita $k$:
\[ h(k) = |\{(i,j) : I(i,j) = k\}| \]
L'istogramma normalizzato e una stima della PDF dell'intensita: $p(k) = h(k) / (M \cdot N)$.
\textbf{2.2 Equalizzazione dell'Istogramma}
Trasformazione che rende l'istogramma il piu possibile uniforme. La funzione di trasformazione e la \textbf{CDF}:
\[ T(k) = \sum_{i=0}^{k} p(i) \]
\begin{itemize}
\item \textbf{Effetto}: aumenta il contrasto globale dell'immagine
\item \textbf{Limite}: non preserva il contrasto locale; puo amplificare il rumore
\end{itemize}
\textbf{2.3 Rumore}
Il rumore e una degradazione stocastica dell'immagine.
Modelli principali:
| Tipo | Modello | Caratteristica |
|------|---------|----------------|
| \textbf{Gaussiano additivo} | $I_{noisy} = I + n$, $n \sim \mathcal{N}(0, \sigma^2)$ | Indipendente dal segnale, sempre presente |
| \textbf{Salt \& Pepper} | Pixel casuali a 0 o 255 | Impulsivo, colpisce una frazione di pixel |
| \textbf{Poisson} | $I_{noisy} \sim \text{Poisson}(I)$ | Correlato all'intensita (shot noise) |
| \textbf{Speckle} | $I_{noisy} = I \cdot n$, $n \sim \mathcal{N}(1, \sigma^2)$ | Moltiplicativo, tipico di radar/ultrasuoni |
\medskip\noindent\rule{\textwidth}{0.4pt}\medskip
\end{frame}

\begin{frame}[allowframebreaks]{3. Filtri e Convoluzione}
\begin{terminiblock}LSIS = sistema lineare shift-invariante (stessa risposta ovunque) · PSF = risposta impulsiva · convoluzione = $y=x*h$ · padding = estensione bordi · separable filter = 2D scomposto in 1D (più veloce) · gaussiano = sfocatura a campana · mediano = valore centrale dopo ordinamento · bilaterale = gaussiano + differenza intensità\end{terminiblock}
\textit{ripulire le immagini dal rumore. E il gesto base di qualsiasi elaborazione: come aggiustare una foto sfocata o piena di puntini bianchi prima di usarla.}
\textbf{Esempio}: selfie con buccia d'arancia: filtro mediano rimuove i puntini senza sfocare. Bilaterale liscia la pelle ma lascia nitidi occhi e bocca.
\begin{spiegazione}La \textbf{convoluzione} è il mattone fondamentale: $y = x * h$, dove $h$ è il kernel (o PSF — Point Spread Function) che scorre sull'immagine come un timbro, moltiplicando i suoi pesi per i pixel sottostanti e sommando. Il sistema è \textbf{LSIS} — Lineare (sovrapposizione: input somma → output somma) e Shift-Invariante (stessa operazione in ogni punto). Ai bordi serve il \textbf{padding} (zero, replicate, reflect) per gestire i pixel fuori immagine.\end{spiegazione}

\textbf{Filtri lineari classici}: . Il \textbf{box filter} fa la media semplice: kernel $k \times k$ con tutti $1/k^2$. Effetto: sfocatura (low-pass), ma sfoca ANCHE i bordi. Il \textbf{gaussiano} $G(x,y)=\frac{1}{2\pi\sigma^2}e^{-\frac{x^2+y^2}{2\sigma^2}}$ è più elegante: pesi a campana — il pixel centrale conta di più, quelli lontani sempre meno. $\sigma$ controlla l'intensità dello sfocamento. Il \textbf{vantaggio cruciale}: è \textbf{separabile}, $G(x,y)=G(x)G(y)$. Invece di un kernel 2D $k \times k$ (costo $O(k^2)$ per pixel), applichi due kernel 1D in sequenza: prima sulle righe, poi sulle colonne — costo $O(k)$. Per un kernel 15×15 passi da 225 a 30 operazioni per pixel.

\textbf{Filtri non lineari}: . Il \textbf{mediano} ordina i pixel nella finestra e prende il valore centrale. Non è esprimibile come convoluzione. Eccelle contro il sale\&pepe: un pixel impulsivo (0 o 255) finisce all'estremità della sequenza ordinata e la mediana lo scarta. \textbf{Preserva i bordi}: sul confine, sceglie il valore del lato con più rappresentanti, invece di mediare e sfocare.

Il \textbf{bilaterale} è il compromesso perfetto tra lisciatura e preservazione dei bordi. Usa DUE pesi moltiplicati: $w = w_s \cdot w_r$. Il peso \textbf{spaziale} $w_s$ è un gaussiano sulla distanza (pixel vicini contano di più). Il peso di \textbf{range} $w_r$ è un gaussiano sulla differenza di intensità (pixel di colore simile contano di più). In una zona uniforme, $w_r \approx 1$ per tutti → si comporta come un gaussiano (liscia). Lungo un bordo, i pixel dall'altra parte hanno intensità molto diversa → $w_r \approx 0$ → NON vengono mediati → il bordo resta nitido. È come colorare dentro le linee: il filtro "capisce" dove sono i confini e non sconfina.
\begin{orale}LSIS = $y=x*h$, stesso filtro ovunque. Mediano = NON sfoca (a differenza del gaussiano). Bilaterale = gaussiano + differenza intensità: preserva bordi. Separabile: 2D = 1D×1D (costo $O(k^2) \to O(k)$). Domanda: «filtro che toglie rumore senza sfocare?» = mediano o bilaterale.\end{orale}
\medskip\noindent\rule{\textwidth}{0.4pt}\medskip
\textbf{3.1 Sistemi LSIS (Lineari Shift-Invarianti)}
Un sistema LSIS e completamente caratterizzato dalla sua \textbf{risposta impulsiva} $h$. L'uscita e la \textbf{convoluzione} dell'ingresso con $h$:
\[ y = x * h \]
Proprieta: commutativa, associativa, distributiva.
\textbf{3.2 Convoluzione 1D e 2D}
\begin{itemize}
\item Il kernel $H$ viene "ribaltato" e fatto scorrere sull'immagine
\item Ai bordi si usa \textbf{padding} (zero, replicate, reflect)
\end{itemize}
\textbf{3.3 Filtri Lineari}
\textbf{Box Filter} (media uniforme):
\[ H = \frac{1}{k^2} \begin{bmatrix} 1 & \cdots & 1 \\ \vdots & \ddots & \vdots \\ 1 & \cdots & 1 \end{bmatrix} \]
\begin{itemize}
\item Effetto: sfocatura (low-pass), rimuove alte frequenze
\item Costo: $O(k^2)$ per pixel
\end{itemize}
\textbf{Filtro Gaussiano}:
\[ G(x,y) = \frac{1}{2\pi\sigma^2} e^{-\frac{x^2+y^2}{2\sigma^2}} \]
\begin{itemize}
\item Sfocatura piu naturale del box filter (pesi maggiori al centro, minori ai bordi)
\item \textbf{Separabile}: $G(x,y) = G(x) \cdot G(y)$ -> costo $O(k)$ per asse invece di $O(k^2)$
\item $\sigma$ controlla l'ampiezza: $\sigma$ grande = piu sfocatura
\end{itemize}
\textbf{3.4 Filtri Non Lineari}
\textbf{Filtro Mediana}:
\begin{itemize}
\item Sostituisce ogni pixel con la mediana del vicinato
\item \textbf{Non lineare} (non esprimibile come convoluzione)
\item Eccellente per rimuovere \textbf{salt \& pepper} preservando i bordi
\end{itemize}
\textbf{Filtro Bilaterale}:
\[ I_{filtered}[i,j] = \frac{\sum_{(m,n) \in \Omega} w_s(m,n) \cdot w_r(I[m,n], I[i,j]) \cdot I[m,n]}{\sum_{(m,n) \in \Omega} w_s(m,n) \cdot w_r(I[m,n], I[i,j])} \]
Due pesi:
\begin{itemize}
\item $w_s$: peso spaziale (vicinanza, gaussiano)  pixel vicini pesano di piu
\item $w_r$: peso di range (similarita di intensita, gaussiano)  pixel con intensita simile pesano di piu
\end{itemize}
\medskip\noindent\rule{\textwidth}{0.4pt}\medskip
\end{frame}

\begin{frame}[allowframebreaks]{4. Trasformata di Fourier}
\begin{terminiblock}DFT = da spazio a frequenze · FFT = algoritmo veloce $O(N \log N)$ · magnitudine = ampiezza frequenze · fase = posizione/struttura (più importante) · teo. convoluzione = spazio $*$ freq $\cdot$ · aliasing = frequenze oltre Nyquist si ripiegano · ringing = artefatto da troncamento brusco\end{terminiblock}
\textit{vedere l'immagine "dall'altro lato" — non pixel ma frequenze. Serve a comprimere (JPEG), rimuovere rumori periodici o capire perche un'immagine ha artefatti.}
\textbf{Esempio}: foto di un prato: DFT mostra tanti puntini lontani dal centro (erba, alte frequenze). Foto di un cielo: tutto concentrato al centro (basse frequenze).
\begin{spiegazione}La \textbf{DFT} (Discrete Fourier Transform) è la lente che ti fa vedere un'immagine non come pixel ma come somma di \textbf{sinusoidi}. È come scomporre un accordo musicale nelle note che lo compongono: ogni "nota" è una sinusoide con una certa frequenza. In 2D: $X[u,v] = \sum_m \sum_n I[m,n] e^{-i2\pi(um/M+vn/N)}$. Le \textbf{basse frequenze} (centro dello spettro dopo `fftshift`) corrispondono a zone uniformi — un cielo, una parete. Le \textbf{alte frequenze} (bordi dello spettro) sono dettagli fini — texture d'erba, contorni, rumore. La \textbf{FFT} rende tutto pratico: invece di $O(N^2)$ operazioni, ne bastano $O(N \log N)$ — su un 1024×1024 passi da \textasciitilde{}10\^{}12 a \textasciitilde{}10\^{}7.\end{spiegazione}

Il \textbf{teorema della convoluzione} è il superpotere del filtraggio in frequenza: $x * h \xleftrightarrow{\mathcal{F}} X \cdot H$. Invece di una convoluzione costosa nello spazio, fai due FFT + un prodotto punto-a-punto. Soprattutto, nel dominio della frequenza \textbf{progettare filtri è intuitivo}: un passa-basso è un cerchio che azzera le frequenze alte; un passa-alto azzera quelle basse. Attenzione però ai tagli netti: un \textbf{filtro ideale} (rettangolare in frequenza) produce \textbf{ringing} — artefatti a "eco" visibili come bordi fantasma (lo vedi nelle JPEG troppo compresse).

\textbf{Fase vs Magnitudine}: — il concetto più controintuitivo. La \textbf{magnitudine} $|X|$ dice QUANTA energia c'è a ogni frequenza (contrasto). La \textbf{fase} $\phi$ dice DOVE stanno le strutture (bordi, posizioni). L'esperimento chiave: prendi la fase dell'immagine A e la magnitudine di B → il risultato è riconoscibile come A. Inverti → riconosci B. \textbf{La fase contiene l'informazione strutturale}: se la alteri, l'immagine diventa irriconoscibile anche con la magnitudine perfetta.

\textbf{Aliasing in frequenza}: componenti oltre $f_{Nyquist}$ si ripiegano su frequenze più basse — è lo stesso fenomeno delle ruote del carro, visto dal lato dello spettro. \textbf{Hybrid images}: applicazione spettacolare — combini basse frequenze di un'immagine (volto serio) con alte frequenze di un'altra (stesso volto che sorride). Da vicino vedi il sorriso (alte freq), da lontano l'espressione seria (basse freq). Funziona perché l'occhio umano ha sensibilità diversa alle frequenze in base alla distanza.
\begin{orale}Teorema convoluzione: filtrare nello spazio = moltiplicare in frequenza (velocissimo). FASE più importante della magnitudine (struttura dell'immagine). Aliasing: frequenze > Nyquist si ripiegano. DFT per analisi, FFT per calcolo. Domanda: «fase o magnitudine?» = fase, contiene la struttura.\end{orale}
\medskip\noindent\rule{\textwidth}{0.4pt}\medskip
\textbf{4.1 Numeri Complessi}
Un numero complesso $z = a + ib = r e^{i\theta}$:
\begin{itemize}
\item \textbf{Magnitudine}: $|z| = \sqrt{a^2 + b^2}$ ? "quanto"
\item \textbf{Fase}: $\theta = \arctan(b/a)$ ? "dove"
\end{itemize}
\textbf{4.2 Serie di Fourier}
Ogni segnale periodico $f(t)$ con periodo $T$ si puo scrivere come somma di sinusoidi:
\[ f(t) = \sum_{n=-\infty}^{\infty} c_n e^{i 2\pi n t / T} \]
\textbf{4.3 DFT (Discrete Fourier Transform)}
\textbf{1D}: $X[k] = \sum_{n=0}^{N-1} x[n] e^{-i 2\pi kn/N}$
\textbf{2D}: $X[u,v] = \sum_{m=0}^{M-1}\sum_{n=0}^{N-1} I[m,n] e^{-i 2\pi (um/M + vn/N)}$
\begin{itemize}
\item \textbf{FFT}: calcola la DFT in $O(N \log N)$ invece di $O(N^2)$
\item Simmetria coniugata: $X[u,v] = X^*[-u,-v]$
\end{itemize}
\textbf{4.4 Teorema della Convoluzione}
\[ x * h \xleftrightarrow{\mathcal{F}} X \cdot H \]
La convoluzione nel dominio spaziale diventa \textbf{prodotto} nel dominio della frequenza (e viceversa).
\textbf{Implicazione pratica}: filtrare in frequenza e $O(N \log N)$ contro $O(N \cdot k^2)$ nello spazio.
\textbf{4.5 Filtraggio in Frequenza}
\begin{itemize}
\item \textbf{Low-pass} (taglia alte frequenze): sfocatura, denoising
\item \textbf{High-pass} (taglia basse frequenze): edge detection, sharpening
\item \textbf{Band-pass}: seleziona una banda di frequenze
\end{itemize}
\textbf{4.6 Importanza di Fase e Magnitudine}
\begin{itemize}
\item \textbf{Fase} ? contiene la \textbf{struttura} (bordi, forme, posizione)
\item \textbf{Magnitudine} ? contiene il \textbf{"quanto"} (contrasto globale, energia)
\end{itemize}
\textbf{4.7 Compressione DFT}
Si possono scartare le componenti di Fourier con magnitudine piccola (sotto una soglia). Poche componenti a bassa frequenza catturano l'essenza dell'immagine.
\textbf{4.8 Hybrid Images}
Combinare basse frequenze di un'immagine con alte frequenze di un'altra:
\begin{itemize}
\item Da vicino si vede l'immagine ad alta frequenza
\item Da lontano si vede l'immagine a bassa frequenza
\end{itemize}
\medskip\noindent\rule{\textwidth}{0.4pt}\medskip
\end{frame}

\begin{frame}[allowframebreaks]{5. Computational Imaging e Problemi Inversi}
\begin{terminiblock}forward model = $y=Ax+e$ · ill-posed = viola esistenza/unicità/stabilità · condizionamento = $\sigma_{max}/\sigma_{min}$ · SVD = $A=U\Sigma V^T$ · $\sigma_i$ = valori singolari · rumore amplificato = $\sigma_i$ piccoli amplificano rumore\end{terminiblock}
\textit{ricostruire un'immagine da misure indirette. TAC, risonanza magnetica, deblurring: senza teoria dei problemi inversi queste applicazioni non esisterebbero.}
\textbf{Esempio}: foto mossa di un gatto (deblurring): $y$ = la foto mossa che hai, $x$ = il gatto nitido che vuoi, $A$ = il movimento della fotocamera.
\begin{spiegazione}Il \textbf{forward model} $\boldsymbol{y} = A\boldsymbol{x} + \boldsymbol{e}$ è il mattone concettuale di tutto il CI. $\boldsymbol{x}$ è l'immagine vera che vuoi (nitida), $A$ è l'operatore fisico di acquisizione (sfocatura, proiezioni CT, sottocampionamento MRI), $\boldsymbol{y}$ sono i dati grezzi che misuri, $\boldsymbol{e}$ è il rumore. Il \textbf{problema diretto} (dato $x$, calcola $y=Ax$) è facile. Il \textbf{problema inverso} (dato $y$, trova $x$) è il nostro incubo quotidiano — e quasi sempre \textbf{ill-posed} nel senso di Hadamard, cioè viola almeno una di tre condizioni: \textbf{esistenza} ($A$ potrebbe non coprire tutte le $y$ possibili, specialmente col rumore), \textbf{unicità} (se $\ker(A) \neq \{0\}$, infinite immagini $x$ producono le stesse misure $y$ — classico della TAC con poche angolazioni), \textbf{stabilità} (piccole perturbazioni nei dati vengono amplificate a dismisura nell'inversione).\end{spiegazione}
La \textbf{SVD} $A = U\Sigma V^T$ svela la struttura nascosta del problema. $U$ e $V$ sono basi ortonormali, $\Sigma = \text{diag}(\sigma_1, \sigma_2, \ldots, \sigma_r)$ contiene i \textbf{valori singolari} in ordine decrescente. La soluzione naive si scrive $\boldsymbol{x}_{naive} = \sum_i \frac{\boldsymbol{u}_i^T \boldsymbol{y}}{\sigma_i} \boldsymbol{v}_i$. Il dramma? Quando $\sigma_i$ è piccolo, il fattore $1/\sigma_i$ esplode e amplifica qualsiasi briciola di rumore lungo la direzione $\boldsymbol{v}_i$. È come un microfono al massimo: senti la voce, ma senti anche ogni fruscio. Le componenti ad alta frequenza dell'immagine sono associate proprio ai $\sigma_i$ più piccoli, quindi il rumore amplificato si manifesta come oscillazioni rapide e artefatti. Il \textbf{numero di condizionamento} $\sigma_{max}/\sigma_{min}$ quantifica la gravità: più è alto, peggio è. Per risolvere, non basta invertire $A$ — serve \textbf{conoscenza a priori} su come deve apparire una soluzione sensata: qui entra la regolarizzazione.
\begin{orale}Hadamard: servono esistenza + unicità + stabilità. Ill-posed se manca una. $\sigma_i$ piccoli dividono per numeri piccoli = amplificano rumore. Domanda tipica: «perché ill-posed?» = $\sigma_i$ piccoli fanno esplodere il rumore.\end{orale}
\medskip\noindent\rule{\textwidth}{0.4pt}\medskip
\textbf{5.1 Il Framework dei Problemi Inversi}
\[ \boldsymbol{y} = A\boldsymbol{x} + \boldsymbol{e} \]
\begin{itemize}
\item $\boldsymbol{x} \in \mathbb{R}^n$: immagine incognita (ground truth)
\item $A \in \mathbb{R}^{m \times n}$: operatore di acquisizione (forward model)
\item $\boldsymbol{y} \in \mathbb{R}^m$: dati misurati
\item $\boldsymbol{e}$: rumore di misura
\end{itemize}
\textbf{Problema diretto}: dato $x$, calcolare $y = Ax$ ? facile, ben posto.
\textbf{Problema inverso}: dato $y$, trovare $x$ ? difficile, spesso mal posto.
\textbf{5.2 Problemi Ill-Posed (Mal Posti)}
Un problema e \textbf{ben posto} (Hadamard) se soddisfa:
\begin{itemize}
\item \textbf{Esistenza}: la soluzione esiste
\item \textbf{Unicita}: la soluzione e unica
\item \textbf{Stabilita}: la soluzione dipende con continuita dai dati
\end{itemize}
I problemi inversi in imaging sono tipicamente \textbf{ill-posed}.
\textbf{5.3 SVD (Singular Value Decomposition)}
\[ A = U \Sigma V^T \]
\begin{itemize}
\item $U, V$: matrici ortogonali (basi ortonormali)
\item $\Sigma = \text{diag}(\sigma_1, \sigma_2, \ldots, \sigma_r)$: valori singolari, $\sigma_1 \geq \sigma_2 \geq \cdots \geq \sigma_r > 0$
\end{itemize}
\textbf{5.4 Soluzione Naive e Amplificazione del Rumore}
La soluzione naive (inversa diretta):
\[ \boldsymbol{x}_{naive} = A^{-1}\boldsymbol{y} = A^{-1}(A\boldsymbol{x} + \boldsymbol{e}) = \boldsymbol{x} + A^{-1}\boldsymbol{e} \]
In termini SVD:
\[ \boldsymbol{x}_{naive} = \sum_{i=1}^{n} \frac{\boldsymbol{u}_i^T \boldsymbol{y}}{\sigma_i} \boldsymbol{v}_i \]
\textbf{Il problema}: quando $\sigma_i$ e piccolo, $1/\sigma_i$ amplifica enormemente il rumore.
\medskip\noindent\rule{\textwidth}{0.4pt}\medskip
\end{frame}

\begin{frame}[allowframebreaks]{6. Regolarizzazione}
\begin{terminiblock}Tikhonov/L2 = $\|x\|^2$ soluzione liscia · TV = $\|\nabla x\|_1$ preserva bordi · L1 = coefficienti nulli (sparsità) · L-curve = curva per scegliere $\lambda$ · $\lambda$ = trade-off fedeltà vs regolarizzazione\end{terminiblock}
\textit{rendere stabile la ricostruzione di un problema inverso. Senza regolarizzazione le immagini sarebbero inutilizzabili. Scegliere L2 o TV cambia il risultato clinico.}
\textbf{Esempio}: deblurring senza regolarizzazione = esce solo rumore. Con L2 = immagine liscia ma senza bordi. Con TV = bordi netti e rumore sparito.
\begin{spiegazione}L'equazione sacra è $\hat{x} = \arg\min_x \|Ax-y\|^2 + \lambda R(x)$. Due forze in competizione. Il \textbf{data fidelity} $\|Ax-y\|^2$ tira verso i dati misurati: se la ricostruzione $x$ non spiega bene $y$, la penalità sale. Il \textbf{regolarizzatore} $R(x)$ tira verso ciò che sappiamo a priori su come deve apparire un'immagine sensata. $\lambda$ è il bilanciere: piccolo = ascolti i dati (rischi overfitting/rumore), grande = ascolti il prior (rischi underfitting/troppo liscio).\end{spiegazione}
Tre regolarizzatori, tre filosofie diverse. \textbf{Tikhonov (L2)}: $R(x)=\|x\|_2^2$ — penalizza pixel grandi, spinge tutto verso zero. Produce soluzioni \textbf{lisce}, senza spigoli, come una palla da bowling. Ottimo per MRI (naturalmente liscia), pessimo per CT ossea (i bordi netti vengono smussati). \textbf{Total Variation (TV)}: $R(x)=\|\nabla x\|_1$ — penalizza la somma dei valori assoluti del gradiente. Trucco: un salto netto (bordo) costa una quantità fissa in TV, non proporzionale all'altezza del salto. Le oscillazioni (rumore) costano MOLTO di più perché ogni micro-fluttuazione si somma. Risultato: preserva i bordi, appiattisce le regioni uniformi. È un \textbf{origami}: pieghe nette (bordi), facce lisce tra le pieghe. \textbf{L1 / Sparsità}: $R(x)=\|x\|_1$ — preferisce soluzioni con tanti zeri. Base del \textbf{compressed sensing}: puoi ricostruire da molte meno misure di quante ne chiederebbe Nyquist, se l'immagine è sparsa in una base opportuna (DCT, wavelet).
Scegliere $\lambda$ è critico. La \textbf{L-curve} è il metodo empirico più usato: plotti $\log\|Ax-y\|^2$ vs $\log R(x)$ per vari $\lambda$. Ottieni una curva a L; il gomito (punto di massima curvatura) è il $\lambda$ ottimale. Alternative: \textbf{discrepanza di Morozov} (scegli $\lambda$ tale che $\|Ax-y\| \approx \sigma$ del rumore) e \textbf{cross-validation}. Il ponte col deep learning: il regolarizzatore artigianale ($\|\nabla x\|_1$, $\|x\|_2^2$) può essere sostituito da un \textbf{prior appreso} dai dati tramite rete neurale — è il passaggio dai metodi classici a quelli moderni.
\begin{orale}$\hat{x} = \arg\min \|Ax-y\|^2 + \lambda R(x)$. L2 = liscio. TV = preserva bordi ($\|\nabla x\|_1$). $\lambda$ grande = troppo liscio, piccolo = rumoroso. L-curve per scegliere $\lambda$. Domanda: «L2 o TV?» = TV se vuoi bordi netti.\end{orale}
\medskip\noindent\rule{\textwidth}{0.4pt}\medskip
\textbf{6.1 Il Framework Model-Based}
\[ \hat{\boldsymbol{x}} = \arg\min_{\boldsymbol{x}} \underbrace{\|A\boldsymbol{x} - \boldsymbol{y}\|_2^2}_{\text{data fidelity}} + \lambda \underbrace{R(\boldsymbol{x})}_{\text{regolarizzazione}} \]
\textbf{6.2 Scelta di $\lambda$}
\begin{itemize}
\item $\lambda \to 0$: overfitting (soluzione rumorosa)
\item $\lambda \to \infty$: underfitting (soluzione troppo liscia)
\item $\lambda$ ottimale: bilancia fedelta e plausibilita
\end{itemize}
\textbf{6.3 Tipi di Regolarizzazione}
| Regolarizzatore | Formula | Effetto |
|-----------------|---------|---------|
| \textbf{Tikhonov} ($L_2$) | $R(x) = \|x\|_2^2$ | Soluzione liscia, penalizza grandi valori |
| \textbf{TV} (Total Variation) | $R(x) = \|\nabla x\|_1$ | Preserva bordi, piecewise constant |
| \textbf{Sparsita} ($L_1$) | $R(x) = \|x\|_1$ | Soluzione sparsa (compressed sensing) |
\textbf{6.4 Collegamento con il Deep Learning}
Il framework model-based classico usa regolarizzatori \textbf{hand-crafted}. Il deep learning sostituisce $R(\boldsymbol{x})$ con un \textbf{prior appreso} dai dati.
\medskip\noindent\rule{\textwidth}{0.4pt}\medskip
\end{frame}

\section{PARTE II - DEEP LEARNING PER IMAGING (Modulo 2)}
\begin{frame}{PARTE II - DEEP LEARNING PER IMAGING (Modulo 2)}
\tableofcontents[currentsection]
\end{frame}

\begin{frame}{}
\medskip\noindent\rule{\textwidth}{0.4pt}\medskip
\end{frame}

\begin{frame}[allowframebreaks]{7. Processing Images per Reti Neurali}
\begin{terminiblock}tensore = $(B,C,H,W)$ PyTorch · normalizzazione = mean/std del dataset (NON 0-1) · Dataset = classe che carica dati · DataLoader = iteratore batch+shuffle+workers · augmentation = trasformazioni casuali\end{terminiblock}
\textit{preparare le immagini per la rete neurale. Dimensioni giuste, valori normalizzati, batch organizzati. Senza preprocessing la rete non impara.}
\textbf{Esempio}: caricare 1000 CT per classificare tumori: le ridimensioni a 224x224, normalizzi, le metti in batch da 32. DataLoader fa tutto con 4 workers in parallelo.
\begin{spiegazione}Le immagini entrano nella rete come \textbf{tensori 4D} $(B, C, H, W)$: batch, canali, altezza, larghezza. Ordine PyTorch, diverso da TensorFlow (NHWC). Un batch di 32 RGB 256×256 è `[32, 3, 256, 256]`.\end{spiegazione}
\textbf{Normalizzazione}: NON è una semplice divisione per 255. La procedura standard è la \textbf{standardizzazione} $x = (x - \mu)/\sigma$, con $\mu$ e $\sigma$ calcolati su TUTTO il training set. Le reti imparano meglio quando ogni feature ha media 0 e varianza 1. Per ImageNet: `mean=[0.485, 0.456, 0.406]`, `std=[0.229, 0.224, 0.225]`. Per GAN e diffusion si normalizza invece a [-1, 1] (Tanh in uscita lo richiede).
\textbf{Dataset e DataLoader}: sono due ruoli distinti. Il \textbf{Dataset} definisce COME caricare un singolo campione (`\_\_getitem\_\_`) e quanti ce ne sono (`\_\_len\_\_`). Il \textbf{DataLoader} avvolge il Dataset e orchestra il training: raggruppa in \textbf{batch} per parallelizzare sulla GPU, fa \textbf{shuffle} a ogni epoca (evita correlazioni spurie nell'ordine), e usa \textbf{workers} multipli per precaricare da disco mentre la GPU elabora — altrimenti la GPU resta in idle.
\textbf{Data Augmentation}: applica trasformazioni casuali durante il training (flip, rotazioni, crop, variazioni di luminosità). Aumenta artificialmente la varietà del dataset senza raccogliere nuovi dati → migliore generalizzazione. In imaging medico: attenzione! No flip verticale per raggi X (su/giù ha significato anatomico).
\textbf{Pipeline di ricostruzione}: . \textbf{End-to-end}: $\boldsymbol{y}^\delta \xrightarrow{f_\Theta} \boldsymbol{x}_{pred}$ — la rete impara direttamente la mappa dati grezzi → immagine. \textbf{Ibrida}: $\boldsymbol{y}^\delta \xrightarrow{\text{FBP}} \tilde{\boldsymbol{x}} \xrightarrow{f_\Theta} \boldsymbol{x}_{pred}$ — un metodo classico (FBP) fa una prima ricostruzione grezza, la rete rimuove solo gli artefatti residui. L'ibrido è preferibile: separa fisica (FBP) da apprendimento (rete), richiede meno dati.
\begin{orale}Tensore = (B, C, H, W). NON normalizzare a [0,1]: usa mean/std del dataset (es. ImageNet: mean=[0.485,0.456,0.406]). DataLoader = batch + shuffle + workers paralleli. Augmentation = più dati senza raccoglierli. Domanda: «normalizzazione?» = sottrarre media, dividere std, NON 0-1.\end{orale}
\medskip\noindent\rule{\textwidth}{0.4pt}\medskip
\textbf{7.1 Tensori}
Le immagini per le NN sono \textbf{tensori 4D}: $(B, C, H, W)$
\begin{itemize}
\item $B$: batch size
\item $C$: canali (1=grigio, 3=RGB)
\item $H, W$: altezza e larghezza
\end{itemize}
\end{frame}

\begin{lstlisting}[style=pp]
# Esempio: batch di 32 immagini RGB 256x256
x.shape  # torch.Size([32, 3, 256, 256])
\end{lstlisting}
\begin{frame}{}
\textbf{7.2 Normalizzazione}
\begin{itemize}
\item \textbf{Range [0, 1]}: `x = x.float() / 255.0` o `transforms.ToTensor()`
\item \textbf{Range [-1, 1]}: `x = (x - 0.5) / 0.5` (usato con GAN, diffusion)
\item \textbf{Standardizzazione}: `x = (x - mean) / std` (per canale)
\end{itemize}
\textbf{dtype}: `float32` per il training (bilancio precisione/memoria).
\textbf{7.3 Dataset e DataLoader}
\end{frame}

\begin{lstlisting}[style=pp]
class MayoDataset(Dataset):
    def __init__(self, data_path):
        self.files = sorted(glob(f'{data_path}/*/*.png'))
        self.transform = transforms.Compose([
            transforms.ToTensor(),
            transforms.Resize((256, 256)),
        ])
    def __len__(self): return len(self.files)
    def __getitem__(self, idx):
        return self.transform(Image.open(self.files[idx]).convert('L'))

loader = DataLoader(dataset, batch_size=32, shuffle=True)
\end{lstlisting}
\begin{frame}{}
\textbf{Mayo Dataset}: 3305 immagini train, 327 test - slice CT addominali.
\textbf{7.4 Pipeline di Ricostruzione}
\textbf{End-to-end}: $y^\delta \xrightarrow{f_\Theta} x_{pred}$
\textbf{Hybrid}: $y^\delta \xrightarrow{\text{FBP}} \tilde{x} \xrightarrow{f_\Theta} x_{pred}$
\medskip\noindent\rule{\textwidth}{0.4pt}\medskip
\end{frame}

\begin{frame}[allowframebreaks]{8. PyTorch Essentials}
\begin{terminiblock}autograd = calcolo automatico gradienti · computational graph = grafo forward · backward() = chain rule all'indietro · SGD = $w = w - lr \cdot \nabla w$ · Adam = lr adattivo per parametro · scheduler = riduce lr durante training\end{terminiblock}
\textit{implementare reti neurali con gradienti automatici, GPU e ottimizzatori. È il toolkit essenziale per qualsiasi progetto di deep learning in imaging.}
\textbf{Esempio}: definisci una rete (nn.Conv2d, nn.ReLU), fai forward, chiami loss.backward(), optimizer.step(). In 15 righe una CNN funzionante.
\begin{spiegazione}\textbf{Autograd} è il motore di differenziazione automatica di PyTorch. Quando operi su tensori con `requires\_grad=True`, PyTorch costruisce un \textbf{computational graph} (DAG: nodi = operazioni, archi = tensori). `loss.backward()` percorre il grafo all'indietro applicando la \textbf{chain rule}, calcolando i gradienti di TUTTI i parametri in un colpo solo — come un domino: la loss è l'ultima pedina, cadendo fa cadere tutte le precedenti. Durante \textbf{valutazione} si usa `torch.no\_grad()`: disabilita il grafo, risparmia memoria e accelera.\end{spiegazione}
\textbf{Ottimizzatori}: l'aggiornamento dei pesi. \textbf{SGD} $w = w - lr \cdot \nabla w$ è il più semplice ma oscilla. \textbf{SGD + momentum} aggiunge una media mobile del gradiente: $v = \mu v + \nabla L$, $w = w - lr \cdot v$ — smorza le oscillazioni e accelera nelle direzioni coerenti (come una palla che rotola giù accumulando quantità di moto). \textbf{Adam} (default oggi) combina momentum + \textbf{adaptive learning rate}: scala lr per ogni parametro individualmente usando la media mobile del gradiente al quadrato. Parametri con gradienti piccoli e consistenti ricevono lr più alto; parametri rumorosi, lr più basso. Iperparametri: $\beta_1=0.9$ (momentum), $\beta_2=0.999$ (RMSprop).
Lo \textbf{scheduler} riduce lr durante il training. \textbf{CosineAnnealingLR}: segue una curva coseno da $lr_{iniziale}$ a $lr_{minima}$, poi riparte (restart). All'inizio esplori con lr alto, poi raffini con lr basso. Alternative: \textbf{StepLR} (fattore gamma ogni N epoche), \textbf{ReduceLROnPlateau} (riduce quando la loss smette di migliorare).
\begin{orale}Autograd = grafo delle operazioni. backward() = gradienti con chain rule. `torch.no\_grad()` quando valuti (non serve grafo). Adam = lr adattivo per parametro (media mobile gradiente + gradiente quadro). Domanda: «backprop?» = chain rule sul grafo computazionale.\end{orale}
\medskip\noindent\rule{\textwidth}{0.4pt}\medskip
\textbf{8.1 Tensori}
\end{frame}

\begin{lstlisting}[style=pp]
x = torch.randn(3, 4)
x.shape, x.dtype, x.device
x.requires_grad = True
\end{lstlisting}
\begin{frame}{}
\textbf{8.2 Moduli Parametrizzati}
\end{frame}

\begin{lstlisting}[style=pp]
class MyNet(nn.Module):
    def __init__(self):
        super().__init__()
        self.linear = nn.Linear(784, 10)
    def forward(self, x):
        return self.linear(x)
\end{lstlisting}
\begin{frame}{}
\textbf{8.3 Autograd}
\end{frame}

\begin{lstlisting}[style=pp]
x = torch.tensor([2.0], requires_grad=True)
y = x**2 + 3*x
y.backward()
print(x.grad)  # tensor([7.0]) = 2*2 + 3
\end{lstlisting}
\begin{frame}{}
\textbf{8.4 Ottimizzatori}
| Ottimizzatore | Formula | Note |
|---------------|---------|------|
| \textbf{SGD} | $\theta \leftarrow \theta - \eta \nabla L$ | Semplice, puo oscillare |
| \textbf{SGD + momentum} | $v \leftarrow \mu v + \nabla L$; $\theta \leftarrow \theta - \eta v$ | Riduce oscillazioni |
| \textbf{Adam} | Combina momentum + RMSProp | Default nella pratica |
\textbf{Learning Rate Scheduler}: `CosineAnnealingLR`
\textbf{8.5 IPPy (libreria del corso)}
\end{frame}

\begin{lstlisting}[style=pp]
IPPy/
+-- operators/   # Forward operators (Blurring, Radon, etc.)
+-- solvers/     # Solvers classici (Tikhonov, ISTA, etc.)
+-- nn/          # Architetture neurali (UNet, DiffusionUNet, etc.)
+-- utilities/   # Device detection, noise generation, etc.
\end{lstlisting}
\begin{frame}{}
\medskip\noindent\rule{\textwidth}{0.4pt}\medskip
\end{frame}

\begin{frame}[allowframebreaks]{9. Da Machine Learning a Neural Networks}
\begin{terminiblock}ReLU = $\max(0,x)$ evita vanishing gradient · sigmoid = $1/(1+e^{-x})$ satura · softmax = output in probabilità · MLP = fully-connected · universal approx = 1 hidden layer approssima qualsiasi funzione · feature = rappresentazioni intermedie · overfitting = memorizza training non generalizza\end{terminiblock}
\textit{capire come un neurone artificiale impara: combinazione lineare + attivazione + ottimizzazione. Serve a risolvere problemi complessi come classificazione e segmentation.}
\textbf{Esempio}: riconoscere cane/gatto: MLP prende tutti i pixel, ReLU trova pattern, softmax dice "75\% cane, 25\% gatto".
\begin{spiegazione}\textbf{ReLU} $\max(0,x)$ è l'attivazione di default per i layer nascosti. Vantaggio cruciale: \textbf{non satura}. Sigmoide e tanh per grandi valori positivi/negativi si appiattiscono (derivata \textasciitilde{}0 → \textbf{vanishing gradient}: l'apprendimento si blocca). ReLU ha derivata 1 per input positivi → i gradienti fluiscono attraverso molti layer. Problema: per input negativi la derivata è 0 → neurone \textbf{muore per sempre} (Dying ReLU), non impara più. Soluzione: \textbf{LeakyReLU} $\max(\alpha x, x)$ con $\alpha$ piccolo (es. 0.01) — lascia passare un piccolo gradiente anche per input negativi, mantenendolo vivo.\end{spiegazione}
\textbf{Softmax}: $p_i = e^{x_i} / \sum_j e^{x_j}$: converte logits in probabilità (somma = 1). Standard per classificazione. \textbf{MLP} (Multi-Layer Perceptron): ogni neurone connesso a TUTTI i neuroni del layer successivo (dense layer). \textbf{Universal Approximation Theorem}: un MLP con 1 hidden layer + sufficienti neuroni approssima QUALSIASI funzione continua su un compatto. MA per immagini è impraticabile: 100×100 = 10.000 input × 1000 neuroni hidden = 10 milioni di parametri solo per il primo layer → overfitting garantito.
Le \textbf{feature} hanno struttura \textbf{gerarchica}: layer bassi → bordi e gradienti; layer intermedi → texture, angoli, cerchi; layer alti → parti di oggetti e oggetti completi. Emerge naturalmente dal training, senza programmazione esplicita. Le \textbf{CNN} risolvono il problema dei troppi parametri con il \textbf{weight sharing}: lo stesso kernel 3×3 (solo 9 pesi, non 10.000) scorre su tutte le posizioni, come un timbro. Sfrutta la struttura spaziale: una feature locale (es. un bordo verticale) è utile ovunque nell'immagine.
\begin{orale}ReLU evita vanishing gradient (sigmoide satura), MA muore se input sempre <0 (LeakyReLU risolve). MLP = universal approximator. Domanda: «perché CNN meglio di MLP per immagini?» = weight sharing, meno parametri, sfrutta struttura spaziale.\end{orale}
\medskip\noindent\rule{\textwidth}{0.4pt}\medskip
\textbf{9.1 Limiti dei Modelli Lineari}
Un modello lineare $f(x) = Wx + b$ puo rappresentare solo relazioni lineari.
\textbf{9.2 Funzioni di Attivazione}
| Attivazione | Formula | Uso |
|-------------|---------|-----|
| \textbf{ReLU} | $\max(0, x)$ | Default per hidden layers |
| \textbf{LeakyReLU} | $\max(\alpha x, x)$ | Evita "dying ReLU" |
| \textbf{Sigmoid} | $1/(1+e^{-x})$ | Output binario [0,1] |
| \textbf{Tanh} | $(e^x - e^{-x})/(e^x + e^{-x})$ | Output [-1,1] |
| \textbf{GELU} | $x \cdot \Phi(x)$ | Transformers |
| \textbf{SiLU/Swish} | $x \cdot \sigma(x)$ | Modelli generativi |
| \textbf{Softmax} | $e^{x_i}/\sum e^{x_j}$ | Classificazione |
\textbf{9.3 MLP (Multi-Layer Perceptron)}
\textbf{Teorema di approssimazione universale}: un MLP con un solo hidden layer (sufficientemente largo) puo approssimare qualsiasi funzione continua su un compatto.
\textbf{9.4 Feature Gerarchiche}
\begin{itemize}
\item Layer bassi: bordi, texture, gradienti
\item Layer intermedi: forme, pattern locali
\item Layer alti: oggetti, strutture semantiche
\end{itemize}
\textbf{9.5 Training}
\begin{itemize}
\item \textbf{SGD/Minibatch}: aggiornamento su sottoinsiemi
\item \textbf{Backpropagation}: calcolo efficiente dei gradienti
\item \textbf{Generalizzazione}: performance su dati non visti
\end{itemize}
\medskip\noindent\rule{\textwidth}{0.4pt}\medskip
\end{frame}

\begin{frame}[allowframebreaks]{10. CNN (Convolutional Neural Networks)}
\begin{terminiblock}convoluzione 2D = kernel $k\times k$ scorre sull'immagine · stride = passo (2 = output dimezzato) · padding = bordi extra per mantenere dimensione · pooling = max pooling riduce risoluzione · channel progression = $64 \to 128 \to 256$ · translation equivariance = shift input = shift output\end{terminiblock}
\textit{elaborare immagini in modo efficiente. Invece di connettere ogni pixel a ogni neurone (MLP), le CNN riusano lo stesso pattern su tutta l'immagine. Pochi parametri, grande potenza.}
\textbf{Esempio}: kernel 3x3 rileva bordi verticali: STESSI 9 pesi su TUTTA l'immagine. Ovunque c'è un bordo verticale il neurone si attiva (weight sharing).
\begin{spiegazione}Nelle CNN, la "convoluzione" è tecnicamente \textbf{cross-correlazione} (kernel NON ribaltato), ma si chiama convoluzione per convenzione. L'operazione: kernel $k \times k$ scorre sull'input, a ogni posizione moltiplica i pesi per i pixel sottostanti e somma → produce una \textbf{feature map} che segnala dove quel pattern è presente.\end{spiegazione}
Il \textbf{weight sharing} è la rivoluzione: lo STESSO kernel (9 pesi per un 3×3) viene applicato ovunque. Tre vantaggi esplosivi: (1) parametri \textbf{drasticamente ridotti} — un'immagine 200×200 con MLP (100 neuroni hidden) = 4 milioni di parametri vs 576 per una CNN con 64 kernel 3×3 ($\approx$7000× meno); (2) \textbf{translation equivariance}: se l'input si sposta di un pixel, la feature map si sposta dello stesso pixel — riconosci il pattern ovunque si trovi; (3) \textbf{località}: il kernel vede solo un vicinato, che è la struttura naturale delle immagini.
\textbf{Stride}: = passo del kernel (1 = dimensione invariata con padding, 2 = dimezza). \textbf{Padding}: $p = (k-1)/2$ mantiene la dimensione in uscita. \textbf{Pooling} (max o average): prende gruppi 2×2 e li riduce a un valore — riduce dimensionalità e dà \textbf{invarianza a piccole traslazioni}. \textbf{Channel progression} 64→128→256: compensa il dimezzamento spaziale raddoppiando i canali, preservando la capacità rappresentativa.
\textbf{Stack di kernel 3×3 > kernel 7×7}: tre 3×3 hanno campo recettivo 7×7 ma solo $3 \times 9 = 27$ pesi (contro 49) + tre ReLU intermedie contro una sola → funzione più espressiva. È il principio dietro VGG.
\begin{orale}Weight sharing = \textasciitilde{}7000× meno parametri di MLP. Stack 3×3 > 7×7: meno parametri + più ReLU. Channel progression: 64 → 128 → 256. Domanda: «perché più canali in profondità?» = compensa la riduzione spaziale.\end{orale}
\medskip\noindent\rule{\textwidth}{0.4pt}\medskip
\textbf{10.1 Convoluzione 2D nelle CNN}
\[ y[i,j] = \sum_{m}\sum_{n} x[i+m, j+n] \cdot w[m,n] + b \]
Nota: nelle CNN si usa la \textbf{cross-correlazione} (senza ribaltamento del kernel), non la convoluzione matematica.
\textbf{Padding}: aggiunta di pixel ai bordi per controllare la dimensione dell'output.
\begin{itemize}
\item `padding = (k-1)/2` (con stride=1) -> output stessa dimensione dell'input
\end{itemize}
\textbf{10.2 Architettura CNN}
\end{frame}

\begin{lstlisting}[style=pp]
class SimpleCNN(nn.Module):
    def __init__(self):
        super().__init__()
        self.net = nn.Sequential(
            nn.Conv2d(1, 32, 3, padding=1),  # 1?32 canali
            nn.ReLU(),
            nn.Conv2d(32, 64, 3, padding=1), # 32?64 canali
            nn.ReLU(),
            nn.Conv2d(64, 1, 3, padding=1),  # 64?1 canale (output)
        )
    def forward(self, x):
        return self.net(x)
\end{lstlisting}
\begin{frame}{}
\textbf{Pooling (opzionale, non nella SimpleCNN)}:
\begin{itemize}
\item \textbf{Max Pooling}: seleziona il valore massimo in una finestra $2 \times 2$ ? dimezza risoluzione, mantiene le feature piu attive
\item \textbf{Average Pooling}: media i valori nella finestra -> dimezza risoluzione, smoothing
\item \textbf{Perche si usa il pooling?}: Riduce la risoluzione spaziale (e il costo computazionale) e aumenta il receptive field dei layer successivi. Nelle CNN per ricostruzione (senza downsampling), il pooling non si usa perche si vuole mantenere la risoluzione piena.
\end{itemize}
\textbf{10.3 Translation Equivariance}
Le CNN sono \textbf{equivarianti} alle traslazioni: se l'input si sposta, l'output si sposta nello stesso modo. Questo perche i pesi sono \textbf{condivisi} su tutta l'immagine (weight sharing).
\textbf{10.4 Training End-to-End per Ricostruzione}
Nel corso si usa un approccio \textbf{sintetico}: i dati di training sono generati online:
\[ \boldsymbol{y} = K\boldsymbol{x} + \boldsymbol{e} \]
dove $K$ e un operatore di blurring e $\boldsymbol{e}$ e rumore gaussiano. Non serve un dataset paired.
\textbf{Attivazione finale}:
\begin{itemize}
\item CNN standard: ReLU o nessuna (output $\geq 0$)
\item ResCNN: Tanh (output in $[-1, 1]$ per il residuo)
\end{itemize}
\medskip\noindent\rule{\textwidth}{0.4pt}\medskip
\end{frame}

\begin{frame}[allowframebreaks]{11. Residual Learning e UNet}
\begin{terminiblock}residual learning = impara $f(x)-x$ (più facile di $f(x)$) · receptive field = regione input che influenza un neurone · skip connection = salta layer (autostrada per gradiente) · encoder-decoder = imbuto comprime + ricostruisce · bottleneck = punto di massima compressione\end{terminiblock}
\textit{segmentare e ricostruire immagini mediche. La UNet è lo standard per CT e MRI. Le skip connections hanno reso possibile addestrare reti molto profonde.}
\textbf{Esempio}: segmentare polmoni in CT: UNet comprime ("cos'è?"), skip riporta dettagli spaziali ("dov'è?"), decoder ricostruisce la maschera.
\begin{spiegazione}\textbf{ResNet} introduce il \textbf{residual learning}: invece di $f(x)$, impara il \textbf{residuo} $f(x)-x$. Quando input e output sono simili (ricostruzione, denoising, super-risoluzione), è molto più facile imparare una piccola correzione che l'intera mappa — come ritoccare un disegno anziché rifarlo da zero. Formula: $x_{pred} = y^\delta + f_\Theta(y^\delta)$, dove $f_\Theta(y^\delta) \approx x - y^\delta$. La \textbf{skip connection} che realizza questo pattern è anche un'\textbf{autostrada per il gradiente}: senza skip, il gradiente deve attraversare TUTTI i layer sequenzialmente (vanishing). Con skip, ha una via diretta verso i layer iniziali — ecco perché le ResNet arrivano a 152 layer mentre VGG si fermava a 19.\end{spiegazione}
\textbf{Receptive field}: (campo recettivo): la regione dell'input che influenza un neurone profondo. Formula ricorsiva: $r_L = r_{L-1} + (k_L-1) \prod_{i=1}^{L-1} s_i$. Più la rete è profonda, più il receptive field cresce. Per segmentazione è cruciale coprire l'intera struttura anatomica, non solo dettagli locali.
\textbf{UNet}: è l'architettura regina dell'imaging medico. Struttura \textbf{encoder-decoder} simmetrica: l'\textbf{Encoder} (contraente) → blocchi Conv + MaxPool, dimezza $H \times W$, raddoppia canali, risponde a "COSA c'è?" (contesto semantico). Il \textbf{Bottleneck} → massima compressione, massimo receptive field, rappresentazione più astratta. Il \textbf{Decoder} (espansivo) → blocchi Conv + UpSampling/ConvTranspose, raddoppia $H \times W$, dimezza canali, risponde a "DOV'È?" (risoluzione spaziale). Le \textbf{Skip connections} (qui CONCATENAZIONI, non somme come ResNet) collegano ogni livello encoder al corrispondente decoder → riportano i dettagli spaziali persi nella compressione.
Varianti: \textbf{Residual UNet} (blocchi residuali), \textbf{Attention UNet} (attention gates sulle skip), \textbf{UNet++} (skip connesse a livelli intermedi).
\begin{orale}ResNet: $f(x)-x$ = più facile se input $\approx$ output. UNet: encoder-decoder + skip (dettagli). Bottleneck = compressione. Domanda: «UNet?» = imbuto (contesto) + skip (dettagli).\end{orale}
\medskip\noindent\rule{\textwidth}{0.4pt}\medskip
\textbf{11.1 ResCNN (Residual CNN)}
Invece di imparare la mappa diretta $y^\delta \to x$, la ResCNN impara il \textbf{residuo}:
\[ f_\Theta(y^\delta) \approx x - y^\delta \]
\[ x_{pred} = y^\delta + f_\Theta(y^\delta) \]
\textbf{Vantaggio}: il residuo e tipicamente piccolo e strutturato -> piu facile da imparare.
\textbf{Skip connection}: $x_{pred} = y^\delta + f_\Theta(y^\delta)$ e essa stessa una skip connection globale.
\textbf{11.2 Receptive Field}
Il \textbf{receptive field} di un neurone e la regione dell'input che influenza la sua attivazione.
Per una CNN con $L$ layer convoluzionali:
\[ r_L = r_{L-1} + (k_L - 1) \prod_{i=1}^{L-1} s_i \]
dove $k_L$ e la dimensione del kernel e $s_i$ lo stride al layer $i$.
\textbf{Problema}: CNN poco profonde hanno receptive field piccolo -> non catturano contesto globale. Soluzione nell'UNet: downsampling + bottleneck permette ai layer profondi di avere receptive field molto grande (tutta l'immagine dopo vari stride 2).
\textbf{11.3 UNet}
Architettura \textbf{encoder-decoder} multi-scala con \textbf{skip connections}:
\end{frame}

\begin{lstlisting}[style=pp]
Encoder (downsampling):       Decoder (upsampling):
  x -> Conv -> 64               512 -> Conv -> 256 -- concat(skip) -> Conv
 64 -> Conv -> 128              256 -> Conv -> 128 -- concat(skip) -> Conv
128 -> Conv -> 256              128 -> Conv ->  64 -- concat(skip) -> Conv
256 -> Conv -> 512 (bottleneck)  64 -> Conv -> output
\end{lstlisting}
\begin{frame}{}
\begin{itemize}
\item \textbf{Encoder}: riduce risoluzione (downsampling con stride 2 o pooling), aumenta canali -> cattura contesto globale a varie scale
\item \textbf{Decoder}: aumenta risoluzione (upsampling con trasposta o interpolazione), riduce canali -> ricostruisce dettagli spaziali
\item \textbf{Skip connections (concat)}: collegano encoder e decoder alla stessa risoluzione tramite \textbf{concatenazione} (non somma come ResNet) -> preservano dettagli spaziali che altrimenti andrebbero persi nel bottleneck
\item \textbf{Bottleneck}: rappresentazione compatta al centro con il maggior numero di canali e la minima risoluzione
\end{itemize}
\textbf{Varianti}:
\begin{itemize}
\item \textbf{Residual UNet}: blocchi residui al posto delle conv semplici
\item \textbf{Attention UNet}: attention gates sulle skip connections
\item \textbf{UNet++}: skip connections nidificate e dense
\end{itemize}
\medskip\noindent\rule{\textwidth}{0.4pt}\medskip
\end{frame}

\begin{frame}[allowframebreaks]{12. Vision Transformers e Loss Design}
\begin{terminiblock}patch embedding = dividi img in patch e proietta in vettori · self-attention = $\text{softmax}(QK^T/\sqrt{d})V$ · positional encoding = vettore posizione (attention è invariante) · PSNR = pixel a pixel >30 dB buono · SSIM = struttura locale · LPIPS = feature VGG percettiva\end{terminiblock}
\textit{scegliere l'architettura (Transformer vs CNN) e la loss function giusta. MSE dà risultati sfocati, LPIPS dà risultati realistici. La differenza tra "pixel uguali" e "sembra vero".}
\textbf{Esempio}: MSE ricostruisce un volto: faccia sfocata (media di tutti i volti). LPIPS: volto nitido con occhi e bocca definiti.
\begin{spiegazione}I \textbf{ViT} (Vision Transformer) adattano l'architettura nata per NLP al dominio immagini. Primo passo: \textbf{patch embedding} — l'immagine $H \times W$ viene divisa in patch $P \times P$ (tipico 16×16), ogni patch appiattita e proiettata linearmente in un embedding. Per 224×224 con patch 16×16: $N = 14 \times 14 = 196$ token. Aggiungi un token [CLS] + \textbf{positional encoding} (appreso o sinusoidale) perché la self-attention è permutation-invariant — senza, una mela a sinistra = una mela a destra.\end{spiegazione}
\textbf{Self-attention}: $\text{softmax}(Q K^T / \sqrt{d_k}) V$. Ogni token genera Query (cosa cerco), Key (cosa offro), Value (contenuto). L'attention calcola il prodotto scalare Q·K tra ogni coppia di token, pesato da softmax → ogni token "guarda" TUTTI gli altri e capisce quanto sono correlati. Divisione per $\sqrt{d_k}$: stabilizza i gradienti (evita saturazione softmax). \textbf{MHSA}: $h$ self-attention parallele (es. $h=8$), ognuna impara relazioni diverse (locali, globali, texture). Costo: $O(N^2)$ — il principale svantaggio vs le CNN che sono $O(N)$. I Transformer catturano contesto globale ma pagano in computazione.
\textbf{Loss function}: — determinano COSA la rete ottimizza. \textbf{MSE/L2} $\frac{1}{n}\|x - \hat{x}\|^2$: pixel a pixel. Semplice ma produce immagini \textbf{sfocate} — matematicamente minimizza $E[x|y]$, la media su tutte le ricostruzioni possibili, e la media di immagini diverse è sfocata. \textbf{PSNR}: $10 \log_{10}(MAX^2 / \text{MSE})$ dB. >40 dB = ottimo, 30-40 = buono, <25 = scarso. Metrica più usata in letteratura. \textbf{SSIM}: confronta luminanza, contrasto e struttura in finestre locali. Range [-1,1], 1 = identiche. Cattura la percezione di struttura: due immagini possono avere PSNR alto ma SSIM basso se la struttura differisce. \textbf{LPIPS} (più moderno): calcola la distanza nello spazio di feature di una VGG pre-addestrata: $\text{LPIPS}(x, \hat{x}) = \sum_l \frac{1}{H_l W_l} \|\phi_l(x) - \phi_l(\hat{x})\|^2$. Non chiede "i pixel sono uguali?" ma "sembra la stessa cosa?" — risultati percettivamente molto più realistici.
\begin{orale}Attention = softmax(QK\^{}T/$\surd$d)V. Self-attention = contesto globale ma $O(N^2)$. MSE = media = sfoca. LPIPS = percettivo (feature VGG). SSIM = struttura locale. Domanda: «loss migliore?» = LPIPS per qualità visiva, PSNR per fedeltà pixel.\end{orale}
\medskip\noindent\rule{\textwidth}{0.4pt}\medskip
\textbf{12.1 Vision Transformer (ViT)}
I Transformer, nati per NLP, sono stati adattati alle immagini.
\textbf{Patch Embedding}: l'immagine $H \times W$ viene divisa in patch $P \times P$:
\begin{itemize}
\item Numero di patch: $N = (H/P) \times (W/P)$
\item Ogni patch -> vettore tramite proiezione lineare
\item Sequenza di $N$ token + token [CLS] opzionale
\end{itemize}
\textbf{Self-Attention}:
\[ \text{Attention}(Q, K, V) = \text{softmax}\left(\frac{Q K^T}{\sqrt{d_k}}\right) V \]
\begin{itemize}
\item $Q$ (Query), $K$ (Key), $V$ (Value): proiezioni lineari dell'input
\item $\sqrt{d_k}$: scaling per evitare gradienti troppo piccoli
\item Ogni token "guarda" tutti gli altri token ? \textbf{contesto globale}
\end{itemize}
\textbf{Multi-Head Self-Attention (MHSA)}: $h$ attention heads in parallelo, ognuno impara relazioni diverse.
\textbf{Positional Encoding}: aggiunto ai patch embeddings per dare informazione di posizione (senza di esso, il transformer e permutation-invariant).
\textbf{12.2 Encoder-Decoder per Image-to-Image}
Per tasks di ricostruzione, si usa un ViT encoder + decoder (spesso convoluzionale):
\begin{itemize}
\item Encoder ViT: cattura dipendenze globali
\item Decoder: ricostruisce l'immagine a risoluzione piena
\end{itemize}
\textbf{CNN vs ViT}:
\begin{itemize}
\item CNN: efficiente per feature locali, induttive bias forte (localita, equivarianza)
\item ViT: contesto globale, meno bias induttivo, richiede piu dati
\end{itemize}
\textbf{12.3 Metriche di Qualita}
\textbf{PSNR} (Peak Signal-to-Noise Ratio):
\[ \text{PSNR} = 10 \log_{10}\left(\frac{\text{MAX}^2}{\text{MSE}}\right) \quad [\text{dB}] \]
\begin{itemize}
\item Piu alto = migliore. Tipico: 25-40 dB per imaging medico.
\end{itemize}
\textbf{SSIM} (Structural Similarity Index):
\[ \text{SSIM}(x,y) = \frac{(2\mu_x\mu_y + c_1)(2\sigma_{xy} + c_2)}{(\mu_x^2 + \mu_y^2 + c_1)(\sigma_x^2 + \sigma_y^2 + c_2)} \]
\begin{itemize}
\item Range $[-1, 1]$, 1 = identico. Considera luminanza, contrasto e struttura.
\end{itemize}
\textbf{LPIPS} (Learned Perceptual Image Patch Similarity):
\begin{itemize}
\item Distanza nello spazio feature di una rete pre-addestrata (VGG/AlexNet)
\item Piu basso = piu simile percettivamente
\end{itemize}
\textbf{12.4 Loss Functions}
| Loss | Formula | Caratteristica |
|------|---------|----------------|
| \textbf{MSE} ($L_2$) | $\frac{1}{n}\|x - \hat{x}\|_2^2$ | Media, produce immagini sfocate |
| \textbf{L1} | $\frac{1}{n}\|x - \hat{x}\|_1$ | Mediana, piu sharp di MSE |
| \textbf{Perceptual} | $\|\phi(x) - \phi(\hat{x})\|_2^2$ | Feature space (VGG), dettagli realistici |
\medskip\noindent\rule{\textwidth}{0.4pt}\medskip
\end{frame}

\begin{frame}[allowframebreaks]{13. Problemi Cross-Domain}
\begin{terminiblock}sinogramma = proiezioni a vari angoli · Radon $\mathcal{R}$ = integrale densità lungo linea · FBP = filtro Ramp + backprojection · ramp filter = $|\omega|$ amplifica alte freq · apodization = finestra per ridurre rumore · streaking = strisce con pochi angoli\end{terminiblock}
\textit{fare CT dal mondo reale. Dalla macchina a raggi X all'immagine finale: sinogramma, FBP, e pipeline ibrida con deep learning. Esempio completo di problema inverso risolto.}
\textbf{Esempio}: CT torace: raggi X girano attorno al paziente, misurano attenuazione = sinogramma. FBP filtra + backprojection = immagine finale.
\begin{spiegazione}La CT è il \textbf{problema cross-domain} per eccellenza: il dominio dei dati (sinogramma) è diverso dal dominio dell'immagine (sezione anatomica). Forward model: $y = \mathcal{R}x + e$, dove $\mathcal{R}$ è la \textbf{Trasformata di Radon} — per ogni angolo $\theta$ misura l'integrale di densità lungo linee rette: $\mathcal{R}f(\theta, s) = \int f(s\cos\theta - t\sin\theta, s\sin\theta + t\cos\theta) dt$. I raggi X attraversano il corpo, vengono attenuati dai tessuti, e i rivelatori misurano l'intensità residua. L'insieme di tutte le proiezioni a vari angoli è il \textbf{sinogramma}.\end{spiegazione}
\textbf{FBP}: (Filtered Backprojection), il metodo classico, ha due passaggi: (1) \textbf{Ramp filter} $|\omega|$ — filtra le proiezioni amplificando le alte frequenze per compensare lo smoothing intrinseco della backprojection (che altrimenti produrrebbe un'immagine sfocata $1/r$). Il ramp filter deriva dal \textbf{Teorema della Fetta di Fourier}: la FT 1D di una proiezione a angolo $\theta$ = una fetta della FT 2D dell'immagine nello stesso angolo. (2) \textbf{Backprojection} — ogni proiezione filtrata viene "spalmata" all'indietro lungo la sua direzione; somma di tutte = immagine finale.
\textbf{Streaking}: artefatti a strisce da sotto-campionamento angolare (pochi angoli → ogni angolo lascia una striscia visibile). È l'aliasing nello spazio angolare. \textbf{Apodization}: applica una finestra (Hann, Hamming) al ramp filter per attenuare le alte frequenze → meno rumore ma meno risoluzione (stesso trade-off di Tikhonov). \textbf{Pipeline ibrida}: $y \xrightarrow{\text{FBP}} \tilde{x} \xrightarrow{f_\Theta} x_{pred}$ — FBP fa la prima ricostruzione rapida (con artefatti), una UNet impara a ripulirla. La fisica è gestita analiticamente, la rete impara solo la correzione → meno dati necessari.
\begin{orale}CT = dalle proiezioni (sinogramma) ricostruisci l'immagine = problema inverso. Radon = integrale densità. FBP = Ramp + backprojection. Apodization = toglie rumore ma perde risoluzione. Domanda: «perché CT è problema inverso?» = misuri proiezioni, devi ricostruire l'immagine.\end{orale}
\medskip\noindent\rule{\textwidth}{0.4pt}\medskip
\textbf{13.1 Il Problema}
In CT (Tomografia Computerizzata), i dati misurati sono \textbf{sinogrammi} (proiezioni a vari angoli), non immagini:
\[ \boldsymbol{y} = \mathcal{R}\boldsymbol{x} + \boldsymbol{e} \]
dove $\mathcal{R}$ e la \textbf{Trasformata di Radon}.
Il dominio dei dati (sinogramma) e diverso dal dominio dell'immagine -> problema \textbf{cross-domain}.
\textbf{13.2 FBP (Filtered Back-Projection)}
Metodo classico per ricostruire da sinogrammi:
\begin{itemize}
\item \textbf{Filtraggio}: le proiezioni sono filtrate con un \textbf{ramp filter} (filtro passa-alto $|f|$ in frequenza) per compensare lo "smoothness" della back-projection
\item \textbf{Back-projection}: ogni proiezione filtrata viene "spalmata" all'indietro lungo la stessa direzione di acquisizione, e tutte le direzioni vengono sommate
\end{itemize}
\textbf{FBP e veloce ma produce artefatti a basse dosi} (rumore, streaking).
\textbf{13.3 Pipeline Ibrida}
\[ \boldsymbol{y} \xrightarrow{\text{FBP}} \tilde{\boldsymbol{x}} \xrightarrow{f_\Theta} \boldsymbol{x}_{pred} \]
\begin{itemize}
\item \textbf{Preprocessing}: FBP converte sinogramma -> immagine (rumorosa)
\item \textbf{Post-processing}: UNet rimuove artefatti e rumore
\end{itemize}
\textbf{Vantaggio}: la rete lavora nel dominio immagine (piu semplice), FBP gestisce la fisica.
\medskip\noindent\rule{\textwidth}{0.4pt}\medskip
\end{frame}

\begin{frame}[allowframebreaks]{14. Deep Generative Models: VAE e GAN}
\begin{terminiblock}ELBO = ricostruzione - KL (obiettivo VAE) · KL = divergenza da $\mathcal{N}(0,I)$ · reparameterization = $z = \mu + \sigma\varepsilon$ (per backprop) · discriminator = distingue reali da generate · minimax = G minimizza D massimizza · mode collapse = G produce poca varietà · WGAN = Wasserstein distance\end{terminiblock}
\textit{generare immagini sintetiche realistiche. Data augmentation, prior per problemi inversi, compressione. VAE è la base stabile, GAN dà risultati più nitidi.}
\textbf{Esempio}: VAE genera 100 facce: un po' sfocate ma tutte diverse (stabile). GAN: nitidissime ma ogni tanto esce la stessa faccia (mode collapse).
\begin{spiegazione}\textbf{VAE} (Variational Autoencoder): unisce autoencoder e inferenza variazionale. Flusso: $x \xrightarrow{E} (\mu, \sigma) \to z = \mu + \sigma \varepsilon \xrightarrow{D} \hat{x}$. L'encoder non produce un vettore latente ma una \textbf{distribuzione} $q(z|x) = \mathcal{N}(\mu(x), \sigma(x)^2)$. Il \textbf{reparameterization trick} $z = \mu + \sigma \odot \varepsilon$ con $\varepsilon \sim \mathcal{N}(0, I)$ rende il campionamento differenziabile: il rumore $\varepsilon$ è esterno, $\mu$ e $\sigma$ passano il gradiente — senza, il sampling stocastico bloccherebbe la backprop.\end{spiegazione}
Loss: \textbf{ELBO} = $\mathbb{E}_{q(z|x)}[\log p(x|z)] - \text{KL}(q(z|x) \| p(z))$. Due forze: (1) \textbf{Ricostruzione} (MSE tra $x$ e $\hat{x}$) — il decoder deve essere fedele. (2) \textbf{KL divergence} — forza $q(z|x)$ verso $\mathcal{N}(0, I)$, rendendo lo spazio latente regolare e campionabile. Bilanciamento delicato: KL forte → latente compatto ma ricostruzione peggiore; KL debole → buona ricostruzione ma latente non usabile come generatore. VAE: \textbf{stabile} (converge sempre) ma produce immagini \textbf{sfocate} — l'ELBO ottimizza una media su tutte le possibili ricostruzioni (stesso motivo di MSE).
\textbf{GAN}: gioco \textbf{minimax} a due. \textbf{Generatore} $G$: rumore $z \to$ immagini $G(z)$. \textbf{Discriminatore} $D$: riceve reali ($x$) e finte ($G(z)$), deve distinguerle. Loss: $\min_G \max_D \mathbb{E}[\log D(x)] + \mathbb{E}[\log(1 - D(G(z)))]$. In pratica si usa la \textbf{non-saturating loss} per $G$: $\min_G -\mathbb{E}[\log D(G(z))]$ (gradiente più forte quando $D$ è bravo). GAN produce immagini \textbf{nitide} ma soffre di \textbf{mode collapse}: $G$ impara poche immagini "sicure" che ingannano $D$ invece di coprire tutta la distribuzione.
\textbf{WGAN}: migliora la stabilità con la \textbf{Wasserstein distance}: $W(p_{data}, p_G) = \sup_{\|f\|_L \leq 1} \mathbb{E}[f(x)] - \mathbb{E}[f(G(z))]$. Il \textbf{critic} $f$ (non più discriminatore) assegna un punteggio di qualità, con vincolo 1-Lipschitz imposto via \textbf{gradient penalty} (WGAN-GP). Il gradiente non satura mai → meno mode collapse.
\textbf{Deep Generative Prior (DGP)}: $z^* = \arg\min_z \|K G(z) - y\|^2$ — ottimizzi $z$ con $G$ congelato, cercando l'immagine nella varietà delle immagini realistiche che meglio spiega i dati. Limite: \textbf{representation error} — se $x_{true}$ non sta nella varietà $\{G(z)\}$, la ricostruzione è imperfetta.
\begin{orale}VAE: ELBO = ricostruzione - KL. Stabile ma sfocato. GAN: minimax, nitido ma mode collapse. WGAN = più stabile. Domanda: «VAE o GAN?» = VAE stabile ma sfocato, GAN nitido ma instabile.\end{orale}
\medskip\noindent\rule{\textwidth}{0.4pt}\medskip
\textbf{14.1 Modelli Generativi vs Discriminativi}
\begin{itemize}
\item \textbf{Discriminativo}: impara $p(y|x)$ o $f: x \to y$ (mappatura diretta)
\item \textbf{Generativo}: impara $p_{data}(x)$ (distribuzione delle immagini) -> puo generare nuovi campioni
\end{itemize}
\textbf{Deep Latent Variable Model (DLVM)}:
\[ \boldsymbol{z} \sim p(\boldsymbol{z}), \quad \boldsymbol{x} \sim p_\Theta(\boldsymbol{x}|\boldsymbol{z}) \]
\begin{itemize}
\item $\boldsymbol{z} \in \mathbb{R}^d$ ($d \ll n$): variabile latente (rappresentazione compressa)
\item $p(\boldsymbol{z}) = \mathcal{N}(0, I)$: prior semplice
\item $p_\Theta(\boldsymbol{x}|\boldsymbol{z})$: decoder (rete neurale)
\end{itemize}
\textbf{Latent vector}: descrizione compressa dell'immagine. Non memorizza ogni pixel, ma codifica i fattori principali (anatomia, struttura, texture).
\textbf{14.2 VAE (Variational Autoencoder)}
Un VAE e un AutoEncoder probabilistico:
\[ \boldsymbol{x} \xrightarrow{E_\phi} (\boldsymbol{\mu}_\phi(\boldsymbol{x}), \boldsymbol{\sigma}_\phi(\boldsymbol{x})) \xrightarrow{\text{sampling}} \boldsymbol{z} \xrightarrow{D_\Theta} \hat{\boldsymbol{x}} \]
\textbf{ELBO} (Evidence Lower Bound):
\[ \log p_\Theta(\boldsymbol{x}) \geq \underbrace{\mathbb{E}_{q_\phi(\boldsymbol{z}|\boldsymbol{x})}[\log p_\Theta(\boldsymbol{x}|\boldsymbol{z})]}_{\text{ricostruzione}} - \underbrace{\text{KL}(q_\phi(\boldsymbol{z}|\boldsymbol{x}) \| p(\boldsymbol{z}))}_{\text{regolarizzazione latente}} \]
\begin{itemize}
\item \textbf{Termine di ricostruzione}: il decoder deve ricostruire bene $\boldsymbol{x}$ a partire da $\boldsymbol{z}$
\item \textbf{Termine KL}: il latente $\boldsymbol{z}$ deve essere vicino a $\mathcal{N}(0,I)$
\end{itemize}
\textbf{Reparameterization trick}:
\[ \boldsymbol{z} = \boldsymbol{\mu} + \boldsymbol{\sigma} \odot \boldsymbol{\varepsilon}, \quad \boldsymbol{\varepsilon} \sim \mathcal{N}(0, I) \]
Rende il sampling differenziabile -> backprop possibile.
\textbf{Training in 2 stadi} (nel corso):
\begin{itemize}
\item \textbf{Pretraining AE}: training deterministico (senza KL) -> buona ricostruzione
\item \textbf{Fine-tuning VAE}: attivazione KL con capacity annealing -> spazio latente regolare
\end{itemize}
\textbf{14.3 GAN (Generative Adversarial Network)}
Due reti in competizione:
\begin{itemize}
\item \textbf{Generator} $G_\Theta$: $\boldsymbol{z} \to \boldsymbol{x}_{fake}$ (decoder-like)
\item \textbf{Discriminator} $D_\Psi$: $\boldsymbol{x} \to$ score reale/falso (encoder-like)
\end{itemize}
\textbf{Obiettivo minimax}:
\[ \min_\Theta \max_\Psi \; \mathbb{E}_{\boldsymbol{x}}[\log D_\Psi(\boldsymbol{x})] + \mathbb{E}_{\boldsymbol{z}}[\log(1 - D_\Psi(G_\Theta(\boldsymbol{z})))] \]
\textbf{Non-saturating loss} (piu stabile per il generator):
\[ \min_\Theta -\mathbb{E}_{\boldsymbol{z}}[\log D_\Psi(G_\Theta(\boldsymbol{z}))] \]
\textbf{Problemi delle GAN}:
\begin{itemize}
\item \textbf{Mode collapse}: il generator produce poca varieta
\item \textbf{Training instabile}: equilibrio delicato tra G e D
\item \textbf{WGAN/WGAN-GP}: usano la distanza di Wasserstein + gradient penalty per stabilita
\item \textbf{Spectral normalization}: controlla la norma dei layer del critic
\end{itemize}
\textbf{Varianti importanti}: DCGAN, Conditional GAN, Pix2Pix, CycleGAN, WGAN, WGAN-GP, LSGAN.
\textbf{14.4 Deep Generative Prior (DGP) per Problemi Inversi}
Un generatore pre-addestrato $G$ definisce un \textbf{prior} sulle immagini:
\[ \hat{\boldsymbol{z}} = \arg\min_{\boldsymbol{z}} \|K G(\boldsymbol{z}) - \boldsymbol{y}^\delta\|_2^2 + \lambda \|\boldsymbol{z}\|_2^2 \]
\[ \hat{\boldsymbol{x}} = G(\hat{\boldsymbol{z}}) \]
\begin{itemize}
\item Il generatore e \textbf{congelato}, si ottimizza solo $\boldsymbol{z}$
\item La soluzione e vincolata al \textbf{range} del generatore: $\mathcal{M} = \{G(\boldsymbol{z}) : \boldsymbol{z} \in \mathbb{R}^d\}$
\item Non servono dati paired per la ricostruzione
\end{itemize}
\textbf{Limiti}:
\begin{itemize}
\item \textbf{Representation error}: se $x_{true} \notin \mathcal{M}$, non puo essere recuperato
\item Ottimizzazione non convessa in $\boldsymbol{z}$
\item GAN prior: latente irregolare, ottimizzazione fragile
\item VAE prior: spazio piu regolare, ma immagini piu sfocate
\end{itemize}
\medskip\noindent\rule{\textwidth}{0.4pt}\medskip
\end{frame}

\begin{frame}[allowframebreaks]{15. Diffusion Models}
\begin{terminiblock}DDPM = aggiunge rumore poi denoizza · forward = $\boldsymbol{x}_t = \sqrt{\bar{\alpha}_t}\boldsymbol{x}_0 + \sqrt{1-\bar{\alpha}_t}\boldsymbol{\varepsilon}$ · noise schedule = $\beta_t$ quanto rumore · score matching = modello impara $\nabla \log p(\boldsymbol{x})$ · DDIM = deterministico salta timestep (10× più veloce)\end{terminiblock}
\textit{generare immagini della massima qualità. DALL-E, Stable Diffusion, Midjourney — tutti diffusion models. Imparano a trasformare rumore puro in immagini perfette.}
\textbf{Esempio}: parti da TV statico (rumore puro). 1000 step di denoising = immagine perfetta di un gatto. DDIM: 50 step bastano.
\begin{spiegazione}I \textbf{diffusion models} imparano a invertire un processo di corruzione graduale. \textbf{Forward process} (DDPM): processo markoviano che aggiunge rumore gaussiano in $T$ passi (tipico $T=1000$). Transizione: $q(x_t | x_{t-1}) = \mathcal{N}(\sqrt{1-\beta_t}\, x_{t-1},\, \beta_t I)$, con $\beta_t$ \textbf{noise schedule} (lineare o cosine). La \textbf{formula chiusa} è la chiave dell'efficienza: $x_t = \sqrt{\bar{\alpha}_t}\, x_0 + \sqrt{1 - \bar{\alpha}_t}\, \varepsilon$, dove $\bar{\alpha}_t = \prod_{s=1}^t (1-\beta_s)$ e $\varepsilon \sim \mathcal{N}(0, I)$. Permette di saltare da $x_0$ a QUALSIASI $x_t$ in un colpo solo — fondamentale per il training perché campioni un $t$ casuale e calcoli $x_t$ istantaneamente. $t$ piccolo ($\bar{\alpha}_t \approx 1$) → quasi $x_0$. $t$ grande ($\bar{\alpha}_t \approx 0$) → quasi puro rumore.\end{spiegazione}
\textbf{Training}: loss sorprendentemente semplice: $\mathbb{E}_{x_0, \varepsilon, t}[\|\varepsilon - \varepsilon_\Theta(x_t, t)\|^2]$. Il modello $\varepsilon_\Theta$ (tipicamente una DiffusionUNet) prende $x_t$ e il timestep $t$, e predice il \textbf{rumore aggiunto} $\varepsilon$, non l'immagine pulita. Equivale a fare \textbf{denoising} a ogni livello. Alternativamente, è \textbf{score matching}: $\varepsilon_\Theta(x_t, t) \approx -\nabla_{x_t} \log p(x_t)$ — il gradiente della log-probabilità punta verso zone di maggiore densità.
\textbf{Sampling DDPM}: da $x_T \sim \mathcal{N}(0, I)$, itera all'indietro: $x_{t-1} = \frac{1}{\sqrt{1-\beta_t}}(x_t - \frac{\beta_t}{\sqrt{1-\bar{\alpha}_t}} \varepsilon_\Theta(x_t, t)) + \sigma_t z$. È \textbf{stocastico}: rimuove rumore ma ne reintroduce un po'. Lento: $T$ valutazioni della rete. \textbf{DDIM}: versione \textbf{deterministica} (stesso $x_T$ → stessa immagine) che permette di \textbf{saltare timestep} ($S \ll T$, es. 50 invece di 1000) → 20× più veloce, stesso modello, zero retraining.
Architettura \textbf{DiffusionUNet}: UNet con time embedding sinusoidale, self-attention ai livelli intermedi, Group Normalization (batch size 1), SiLU, EMA dei pesi per sampling stabile.
\begin{orale}Forward = formula chiusa $x_t$ da $x_0$. Training = MSE sul rumore. Sampling = da rumore a immagine. DDPM = stocastico, DDIM = deterministico + veloce. Domanda: «Diffusion in breve?» = aggiungi rumore forward, togli rumore backward.\end{orale}
\medskip\noindent\rule{\textwidth}{0.4pt}\medskip
\textbf{15.1 Idea Centrale}
I diffusion models generano dati imparando a \textbf{invertire un processo di aggiunta graduale di rumore}.
\textbf{15.2 Forward Process (DDPM)}
Sequenza di variabili latenti $\boldsymbol{x}_0, \boldsymbol{x}_1, \ldots, \boldsymbol{x}_T$:
\[ q(\boldsymbol{x}_t | \boldsymbol{x}_{t-1}) = \mathcal{N}(\sqrt{1-\beta_t}\,\boldsymbol{x}_{t-1},\; \beta_t I) \]
dove $\beta_t \in (0,1)$ e il \textbf{noise schedule}.
\textbf{Formula chiusa} (campionamento diretto a qualsiasi $t$):
\[ \boldsymbol{x}_t = \sqrt{\alpha_t}\,\boldsymbol{x}_0 + \sqrt{1-\alpha_t}\,\boldsymbol{\varepsilon}_t, \quad \boldsymbol{\varepsilon}_t \sim \mathcal{N}(0,I) \]
dove $\alpha_t = \prod_{s=1}^{t}(1-\beta_s)$.
\begin{itemize}
\item $t$ piccolo: $\alpha_t \approx 1$, $\sqrt{1-\alpha_t} \approx 0$ ? immagine quasi pulita
\item $t$ grande: $\alpha_t \approx 0$, $\sqrt{1-\alpha_t} \approx 1$ ? quasi puro rumore gaussiano
\item Per $t = T$ (tipicamente $T=1000$): $\alpha_T \approx 0$, $\boldsymbol{x}_T \sim \mathcal{N}(0,I)$
\end{itemize}
\textbf{15.3 Training (Noise Prediction)}
Una singola rete $\boldsymbol{\varepsilon}_\Theta(\boldsymbol{x}_t, t)$ impara a predire il rumore:
\[ \min_\Theta \; \mathbb{E}_{\boldsymbol{x}_0, \boldsymbol{\varepsilon}_t, t}\left[\|\boldsymbol{\varepsilon}_t - \boldsymbol{\varepsilon}_\Theta(\boldsymbol{x}_t, t)\|_2^2\right] \]
\textbf{Stima dell'immagine pulita}:
\[ \hat{\boldsymbol{x}}_0(\boldsymbol{x}_t, t) = \frac{\boldsymbol{x}_t - \sqrt{1-\alpha_t}\,\boldsymbol{\varepsilon}_\Theta(\boldsymbol{x}_t, t)}{\sqrt{\alpha_t}} \]
\textbf{15.4 Architettura del Denoiser}
\textbf{DiffusionUNet}: UNet time-conditioned con:
\begin{itemize}
\item Residual blocks + Group Normalization + SiLU
\item Self-attention a risoluzioni intermedie
\item \textbf{Sinusoidal embedding} per il timestep $t$ (come positional encoding nei transformer)
\end{itemize}
Il time embedding viene proiettato e \textbf{aggiunto} dentro i residual blocks -> ogni blocco sa a quale livello di rumore sta operando.
\textbf{EMA} (Exponential Moving Average): si mantiene una copia "mediata" dei pesi per sampling piu stabile.
\textbf{15.5 DDPM Sampling (Reverse Process)}
Partendo da rumore puro $\boldsymbol{x}_T \sim \mathcal{N}(0,I)$, si applicano $T$ step di denoising:
\[ \boldsymbol{x}_{t-1} = \frac{1}{\sqrt{1-\beta_t}}\left(\boldsymbol{x}_t - \frac{\beta_t}{\sqrt{1-\alpha_t}}\boldsymbol{\varepsilon}_\Theta(\boldsymbol{x}_t, t)\right) + \sigma_t \boldsymbol{z} \]
dove $\boldsymbol{z} \sim \mathcal{N}(0,I)$ e $\sigma_t$ e la varianza del passo.
\textbf{Lento}: richiede $T$ valutazioni della rete (tipicamente $T = 400-1000$). Ogni valutazione e una forward pass della DiffusionUNet.
\textbf{15.6 DDIM Sampling (Deterministico e Veloce)}
\textbf{DDIM} (Denoising Diffusion Implicit Models) modifica il sampling:
\[ \boldsymbol{x}_{s} = \sqrt{\alpha_s}\,\hat{\boldsymbol{x}}_0 + \sqrt{1-\alpha_s}\,\boldsymbol{\varepsilon}_\Theta(\boldsymbol{x}_t, t) \]
\begin{itemize}
\item \textbf{Deterministico}: stesso $\boldsymbol{x}_T$ ? stessa immagine generata
\item \textbf{Accelerato}: si possono saltare timestep (es. 40 step invece di 400)
\item \textbf{Stesso modello}: non serve ri-addestrare
\end{itemize}
\medskip\noindent\rule{\textwidth}{0.4pt}\medskip
\end{frame}

\begin{frame}[allowframebreaks]{16. Diffusion Models per Problemi Inversi}
\begin{terminiblock}posterior sampling = campiona da $p(\boldsymbol{x}|\boldsymbol{y})$ · likelihood = $\|K\boldsymbol{x} - \boldsymbol{y}\|^2$ · DPS = gradiente likelihood corregge sampling ogni step · DiffPIR = alterna denoising + data consistency · score function = $\nabla \log p(\boldsymbol{x})$ direzione per aumentare probabilità\end{terminiblock}
\textit{unire la qualità dei diffusion models con i dati di misura reali. Ottenere ricostruzioni CT/MRI realistiche che rispettano le misure acquisite. È il fronte della ricerca.}
\textbf{Esempio}: CT con poche angolazioni: DPS usa diffusion come "memoria" delle CT vere, il gradiente likelihood guida la ricostruzione a rispettare le misure.
\begin{spiegazione}Il problema inverso in ottica \textbf{Bayesiana}: $p(x | y) \propto p(y | x) \, p(x)$. Il \textbf{prior} $p(x)$ è la distribuzione delle immagini reali, che il diffusion model impara come \textbf{score function} $\nabla \log p(x_t)$ a TUTTI i livelli di rumore. Vantaggio cruciale sui GAN: il diffusion NON comprime in una varietà $\mathcal{M} = \{G(z)\}$ di dimensione $d \ll n$, ma copre \textbf{tutto lo spazio delle immagini} → nessun representation error (se $x_{true}$ non sta in $\mathcal{M}$, il GAN fallisce; il diffusion no). La \textbf{likelihood} $p(y|x) \propto \exp(-\|Kx - y\|^2 / 2\sigma_y^2)$ misura l'accordo con le misure. Il posterior sampling combina entrambe per campionare da $p(x|y)$.\end{spiegazione}
\textbf{DPS}: (Diffusion Posterior Sampling): a ogni step di denoising corregge la traiettoria con un gradiente di likelihood. Formula: $\nabla_{x_t} \log p_t(x_t | y) = \underbrace{\nabla_{x_t} \log p_t(x_t)}_{\text{dal diffusion}} + \underbrace{\nabla_{x_t} \log p(y | x_t)}_{\text{guida della misura}}$. Problema: $p(y|x_t)$ non è trattabile (la likelihood è definita su $x_0$, non su $x_t$). Soluzione: approssima $\hat{x}_0$ da $x_t$ con la formula chiusa del forward, poi calcola il gradiente $\nabla_{x_t} \|K \hat{x}_0(x_t) - y\|^2$. Update: $x_s \approx R_\Theta(x_t) - \eta \nabla_{x_t} \|K \hat{x}_0 - y\|^2$. Pro: nessun riaddestramento, $K$ intercambiabile. Contro: computazionalmente pesante, gradiente approssimato, $\eta$ delicato da tunare.
\textbf{DiffPIR}: (Plug-and-Play): filosofia diversa — alterna due step come un algoritmo \textbf{proximal} classico. (1) \textbf{Denoising} (prior): un passo DDIM standard. (2) \textbf{Data consistency}: $x \leftarrow x - \tau K^T (Kx - y)$ — proiezione di gradiente che forza il rispetto delle misure. È un \textbf{HQS} (Half Quadratic Splitting) dove il denoiser è sostituito da un diffusion model. Più interpretabile e stabile di DPS: ogni step ha un significato chiaro. Connessione diretta con i metodi classici di ottimizzazione.
\begin{orale}DPS = denoising + gradiente $\nabla_x\|Ax-y\|^2$. DiffPIR = denoising alternato a data consistency. Prior diffusion > GAN: copre tutto lo spazio immagini. Domanda: «Diffusione per problemi inversi?» = prior diffusivo + guida della misura.\end{orale}
\medskip\noindent\rule{\textwidth}{0.4pt}\medskip
\textbf{16.1 Punto di Vista Bayesiano}
Dato il modello di misura $\boldsymbol{y}^\delta = K\boldsymbol{x}^\dagger + \boldsymbol{e}$ con rumore gaussiano:
\[ p(\boldsymbol{x}|\boldsymbol{y}^\delta) \propto \underbrace{p(\boldsymbol{y}^\delta|\boldsymbol{x})}_{\text{likelihood}} \cdot \underbrace{p(\boldsymbol{x})}_{\text{prior}} \]
\begin{itemize}
\item \textbf{Likelihood}: $p(\boldsymbol{y}^\delta|\boldsymbol{x}) \propto \exp\left(-\frac{1}{2\sigma_y^2}\|K\boldsymbol{x} - \boldsymbol{y}^\delta\|_2^2\right)$  descrive l'accordo con le misure
\item \textbf{Prior}: il diffusion model fornisce informazione di denoising/score a molti livelli di rumore  descrive la plausibilita dell'immagine
\end{itemize}
Il diffusion model non da una densita $p(\boldsymbol{x})$ in forma chiusa, ma fornisce \textbf{informazione di score} $\nabla_{\boldsymbol{x}_t} \log p_t(\boldsymbol{x}_t)$ che guida la ricostruzione iterativa.
\textbf{16.2 DPS (Diffusion Posterior Sampling)}
\textbf{Idea}: modificare la traiettoria del reverse diffusion con un gradiente di likelihood.
Score della posterior:
\[ \nabla_{\boldsymbol{x}_t} \log p_t(\boldsymbol{x}_t|\boldsymbol{y}^\delta) = \underbrace{\nabla_{\boldsymbol{x}_t} \log p_t(\boldsymbol{x}_t)}_{\text{dal diffusion model}} + \underbrace{\nabla_{\boldsymbol{x}_t} \log p(\boldsymbol{y}^\delta|\boldsymbol{x}_t)}_{\text{dalla likelihood}} \]
Il gradiente di likelihood e approssimato usando $\hat{\boldsymbol{x}}_0$:
\[ \boldsymbol{g}_t \approx \nabla_{\boldsymbol{x}_t} \frac{1}{2\sigma_y^2}\|K\hat{\boldsymbol{x}}_0(\boldsymbol{x}_t, t) - \boldsymbol{y}^\delta\|_2^2 \]
\textbf{Update}:
\[ \boldsymbol{x}_{s} \approx R_\Theta(\boldsymbol{x}_t, t, s) - \eta\,\boldsymbol{g}_t \]
dove $R_\Theta$ e lo step DDIM e $\eta$ e la guidance strength.
\textbf{Pro}: prior diffusionale sempre nel loop, riusabile per diversi operatori $K$ senza riaddestrare.
\textbf{Contro}: computazionalmente pesante (gradiente ad ogni step), $\eta$ da tuningare, gradiente approssimato.
\textbf{16.3 DiffPIR (Diffusion Plug-and-Play Image Restoration)}
\textbf{Filosofia plug-and-play}: alternare denoising (prior) e data consistency (misure).
Ad ogni timestep:
\begin{itemize}
\item \textbf{Prior step}: step DDIM -> $\boldsymbol{x}_{prior}$ (spinto verso il prior)
\item \textbf{Data consistency}: correzione verso le misure:
\end{itemize}
\[ \boldsymbol{x}_{next} = \boldsymbol{x}_{prior} - \tau K^T(K\boldsymbol{x}_{prior} - \boldsymbol{y}^\delta) \]
dove $\tau > 0$ e uno step size.
\textbf{Pro}: modulare, interpretabile, connesso a metodi di ottimizzazione classici (proximal algorithms).
\textbf{Contro}: non e un campionatore posterio esatto, bilanciamento denoising/correction delicato.
\textbf{16.4 Confronto DPS vs DiffPIR}
| | DPS | DiffPIR |
|---|-----|---------|
| \textbf{Filosofia} | Steering della traiettoria posterio | Operator splitting (alternanza) |
| \textbf{Data consistency} | Gradiente di likelihood | Correzione esplicita $K^T(Kx - y)$ |
| \textbf{Interpretabilita} | Meno trasparente | Piu trasparente (proximal) |
| \textbf{Stabilita} | Sensibile a $\eta$ | Piu stabile con $\tau$ |
| \textbf{Costo} | Gradiente via autograd | Solo forward di $K$ e $K^T$ |
\textbf{16.5 Limitazioni Generali}
\begin{itemize}
\item \textbf{Computazionalmente costosi}: molte valutazioni della rete per ricostruzione
\item \textbf{Sensibili al forward model}: se $K$ e sbagliato, la ricostruzione peggiora
\item **Plausibilita $\neq$ correttezza**: immagini realistiche ma potenzialmente inventate
\item \textbf{Distribution shift}: prior addestrato su una popolazione puo fallire su anatomie diverse
\item \textbf{Iperparametri critici}: guidance strength, noise schedule, numero di step
\end{itemize}
\medskip\noindent\rule{\textwidth}{0.4pt}\medskip
\end{frame}

\section{SCHEDE RIASSUNTIVE}
\begin{frame}{SCHEDE RIASSUNTIVE}
\tableofcontents[currentsection]
\end{frame}

\end{document}